\documentclass[UTF8, 12pt, fontset=fandol, oneside]{ctexart}
\usepackage{researchnotes}
\DeclareMathOperator{\dom}{dom}
\DeclareMathOperator{\diag}{diag}
\DeclareMathOperator{\dist}{dist}
\DeclareMathOperator{\epi}{epi}
\DeclareMathOperator{\hypo}{hypo}

\title{第三章\quad 凸函数}
\author{}
\date{}

\begin{document}

\maketitle

\section*{3.1\quad 基本性质和例子}

\subsection*{3.1.1\quad 定义}

函数 $f:\mathbb{R}^n\to\mathbb{R}$ 是凸的，如果 $\dom f$ 是凸集，且对于任意 $x,y\in\dom f$ 和任意 $0\leq\theta\leq 1$，有
\begin{equation}
f(\theta x+(1-\theta)y)\leq \theta f(x)+(1-\theta)f(y).
\tag{3.1}
\end{equation}
从几何意义上看，上述不等式意味着点 $(x,f(x))$ 和 $(y,f(y))$ 之间的线段，即从 $x$ 到 $y$ 的弦，在函数 $f$ 的图像上方（如图 3.1 所示）。称函数 $f$ 是严格凸的，如果式 (3.1) 中的不等式当 $x\neq y$ 以及 $0<\theta<1$ 时严格成立。称函数 $f$ 是凹的，如果函数 $-f$ 是凸的，称函数 $f$ 是严格凹的，如果 $-f$ 严格凸。

对于仿射函数，不等式 (3.1) 总成立，因此所有的仿射函数（包括线性函数）是既凸且凹的。反之，若某个函数是既凸又凹的，则其是仿射函数。

函数是凸的，当且仅当其在与其定义域相交的任何直线上都是凸的。换言之，函数 $f$ 是凸的，当且仅当对于任意 $x\in\dom f$ 和任意向量 $v$，函数 $g(t)=f(x+tv)$ 是凸的（其定义域为 $\{t\mid x+tv\in\dom f\}$）。这个性质非常有用，因为它容许我们通过将函数限制在直线上来判断其是否是凸函数。

对凸函数的分析已经相当地透彻，本书不再继续深入。例如有这样一个简单的结论：凸函数在其定义域相对内部是连续的；它只可能在相对边界上不连续。

\subsection*{3.1.2\quad 扩展值延伸}

通常可以定义凸函数在定义域外的值为 $\infty$，从而将这个凸函数延伸至全空间 $\mathbb{R}^n$。如果 $f$ 是凸函数，我们定义它的扩展值延伸 $\tilde f:\mathbb{R}^n\to\mathbb{R}\cup\{\infty\}$ 如下
\[
\tilde f(x)=
\begin{cases}
f(x), & x\in\dom f,\\
\infty, & x\notin\dom f.
\end{cases}
\]
延伸函数 $\tilde f$ 是定义在全空间 $\mathbb{R}^n$ 上的，取值集合为 $\mathbb{R}\cup\{\infty\}$。我们也可以从延伸函数 $\tilde f$ 的定义中确定原函数 $f$ 的定义域，即 $\dom f=\{x\mid \tilde f(x)<\infty\}$。

这种延伸可以简化符号描述，此时我们就不需要明确描述定义域或者每次提到 $f(x)$ 时都限定“对于所有的 $x\in\dom f$”。以基本不等式 (3.1) 为例，对于延伸函数 $\tilde f$，可以描述为：对于任意 $x$ 和 $y$，以及 $0<\theta<1$，有
\[
\tilde f(\theta x+(1-\theta)y)\leq \theta \tilde f(x)+(1-\theta)\tilde f(y).
\]
（当 $\theta=0$ 或 $\theta=1$ 时不等式总成立。）当然此时我们应当利用扩展运算和序来理解这个不等式。若 $x$ 和 $y$ 都在 $\dom f$ 内，上述不等式即为不等式 (3.1)；如果有任何一个在 $\dom f$ 外，上述不等式的右端为 $\infty$，不等式仍然成立。再看一个这种表示的例子，设 $f_1$ 和 $f_2$ 是 $\mathbb{R}^n$ 上的两个凸函数。逐点和函数 $f=f_1+f_2$ 的定义域为 $\dom f=\dom f_1\cap\dom f_2$，对于任意 $x\in\dom f$，有 $f(x)=f_1(x)+f_2(x)$。利用扩展值延伸我们可以简单地描述为，对于任意 $x$，$\tilde f(x)=\tilde f_1(x)+\tilde f_2(x)$。在这个方程里，函数 $f$ 的定义域被自动定义为 $\dom f=\dom f_1\cap\dom f_2$，因为当 $x\notin\dom f_1$ 或者 $x\notin\dom f_2$ 时，$\tilde f(x)=\infty$。在此例中，我们就可以利用扩展运算来自动定义定义域。

在不会造成歧义的情况下，本书将用同样的符号来表示一个凸函数及其延伸函数。即假设所有的凸函数都隐含地被延伸了，也就是在定义域外都被定义为 $\infty$。

\noindent\textbf{例 3.1\quad 凸集的示性函数。}\quad 设 $C\subseteq\mathbb{R}^n$ 是一个凸集，考虑（凸）函数 $I_C$，其定义域为 $C$，对于所有的 $x\in C$，有 $I_C(x)=0$。换言之，此函数在集合 $C$ 上一直为零。其扩展值延伸可以描述如下
\[
\tilde I_C(x)=
\begin{cases}
0, & x\in C,\\
\infty, & x\notin C.
\end{cases}
\]
凸函数 $\tilde I_C$ 被称作集合 $C$ 的示性函数。

利用示性函数 $\tilde I_C$，可以更加灵活地定义符号描述。例如，对于在集合 $C$ 上极小化函数 $f$（假设其定义在整个 $\mathbb{R}^n$ 空间）的问题，我们可以给出等价的问题，即在 $\mathbb{R}^n$ 上极小化函数 $f+\tilde I_C$。事实上，函数 $f+\tilde I_C$（按照我们的约定）等价于定义在集合 $C$ 上的函数 $f$。

类似地，可以通过定义凹函数在定义域外都为 $-\infty$ 对其进行延伸。

\subsection*{3.1.3\quad 一阶条件}

假设 $f$ 可微（即其梯度 $\nabla f$ 在开集 $\dom f$ 内处处存在），则函数 $f$ 是凸函数的充要条件是 $\dom f$ 是凸集且对于任意 $x,y\in\dom f$，下式成立
\begin{equation}
f(y)\geq f(x)+\nabla f(x)^T(y-x).
\tag{3.2}
\end{equation}
图 3.2 描述了上述不等式。

由 $f(x)+\nabla f(x)^T(y-x)$ 得出的仿射函数 $y$ 即为函数 $f$ 在点 $x$ 附近的 Taylor 近似。不等式 (3.2) 表明，对于一个凸函数，其一阶 Taylor 近似实质上是原函数的一个全局下估计。反之，如果某个函数的一阶 Taylor 近似总是其全局下估计，那么这个函数是凸的。

不等式 (3.2) 说明从一个凸函数的局部信息（即它在某点的函数值及导数），我们可以得到一些全局信息（如它的全局下估计）。这也许是凸函数的最重要的信息，由此可以解释凸函数以及凸优化问题的一些非常重要的性质。下面是一个简单的例子，由不等式 (3.2) 可以知道，如果 $\nabla f(x)=0$，那么对于所有的 $y\in\dom f$，存在 $f(y)\geq f(x)$，即 $x$ 是函数 $f$ 的全局极小点。

严格凸性同样可以由一阶条件刻画：函数 $f$ 严格凸的充要条件是 $\dom f$ 是凸集且对于任意 $x,y\in\dom f$，$x\neq y$，有
\begin{equation}
f(y)>f(x)+\nabla f(x)^T(y-x).
\tag{3.3}
\end{equation}
对于凹函数，亦存在与之对应的一阶条件：函数 $f$ 是凹函数的充要条件是 $\dom f$ 是凸集且对于任意 $x,y\in\dom f$，下式成立
\[
f(y)\leq f(x)+\nabla f(x)^T(y-x).
\]

\noindent\textbf{一阶凸性条件的证明}

为了证明式 (3.2)，先考虑 $n=1$ 的情况：我们证明可微函数 $f:\mathbb{R}\to\mathbb{R}$ 是凸函数的充要条件是对于 $\dom f$ 内的任意 $x$ 和 $y$，有
\begin{equation}
f(y)\geq f(x)+f'(x)(y-x).
\tag{3.4}
\end{equation}
首先假设 $f$ 是凸函数，且 $x,y\in\dom f$。因为 $\dom f$ 是凸集（某个区间），对于任意 $0<t\leq 1$，我们有 $x+t(y-x)\in\dom f$，由函数 $f$ 的凸性可得
\[
f(x+t(y-x))\leq (1-t)f(x)+tf(y).
\]
将上式两端同除 $t$ 可得
\[
f(y)\geq f(x)+\frac{f(x+t(y-x))-f(x)}{t},
\]
令 $t\to0$，可以得到不等式 (3.4)。

为了证明充分性，假设对 $\dom f$（某个区间）内的任意 $x$ 和 $y$，函数满足不等式 (3.4)。选择任意 $x\neq y$，$0\leq\theta\leq 1$，令 $z=\theta x+(1-\theta)y$。两次应用不等式 (3.4) 可得
\[
f(x)\geq f(z)+f'(z)(x-z),\qquad f(y)\geq f(z)+f'(z)(y-z).
\]
将第一个不等式乘以 $\theta$，第二个不等式乘以 $1-\theta$，并将二者相加可得
\[
\theta f(x)+(1-\theta)f(y)\geq f(z),
\]
从而说明了函数 $f$ 是凸的。

现在来证明一般情况，即 $f:\mathbb{R}^n\to\mathbb{R}$。设 $x,y\in\mathbb{R}^n$，考虑过这两点的直线上的函数 $f$，即函数 $g(t)=f(ty+(1-t)x)$，此函数对 $t$ 求导可得 $g'(t)=\nabla f(ty+(1-t)x)^T(y-x)$。

首先假设函数 $f$ 是凸的，则函数 $g$ 是凸的，由前面的讨论可得 $g(1)\geq g(0)+g'(0)$，即
\[
f(y)\geq f(x)+\nabla f(x)^T(y-x).
\]
再假设此不等式对于任意 $x$ 和 $y$ 均成立，因此若 $ty+(1-t)x\in\dom f$ 以及 $\tilde t y+(1-\tilde t)x\in\dom f$，我们有
\[
f(ty+(1-t)x)\geq f(\tilde t y+(1-\tilde t)x)+\nabla f(\tilde t y+(1-\tilde t)x)^T(y-x)(t-\tilde t),
\]
即 $g(t)\geq g(\tilde t)+g'(\tilde t)(t-\tilde t)$，说明了函数 $g$ 是凸的。

\subsection*{3.1.4\quad 二阶条件}

现在假设函数 $f$ 二阶可微，即对于开集 $\dom f$ 内的任意一点，它的 Hessian 矩阵或者二阶导数 $\nabla^2 f$ 存在，则函数 $f$ 是凸函数的充要条件是，其 Hessian 矩阵是半正定阵：即对于所有的 $x\in\dom f$，有
\[
\nabla^2 f(x)\succeq 0.
\]

对于 $\mathbb{R}$ 上的函数，上式可以简化为一个简单的条件：$f''(x)\geq0$（$\dom f$ 是凸的，即一个区间），此条件说明函数 $f$ 的导数是非减的。条件 $\nabla^2 f(x)\succeq0$ 从几何上可以理解为函数图像在点 $x$ 处具有正（向上）的曲率。关于二阶条件的证明作为习题留给读者完成（习题 3.8）。

类似地，函数 $f$ 是凹函数的充要条件是，$\dom f$ 是凸集且对于任意 $x\in\dom f$，$\nabla^2 f(x)\preceq0$。严格凸的条件可以部分由二阶条件刻画。如果对于任意的 $x\in\dom f$ 有 $\nabla^2 f(x)\succ0$，则函数 $f$ 严格凸。反过来则不一定成立：例如，函数 $f:\mathbb{R}\to\mathbb{R}$，其表达式为 $f(x)=x^4$，它是严格凸的，但是在 $x=0$ 处，二阶导数为零。

\noindent\textbf{例 3.2\quad 二次函数。}\quad 考虑二次函数 $f:\mathbb{R}^n\to\mathbb{R}$，其定义域为 $\dom f=\mathbb{R}^n$，其表达式为
\[
f(x)=(1/2)x^TPx+q^Tx+r,
\]
其中 $P\in\mathbb{S}^n$，$q\in\mathbb{R}^n$，$r\in\mathbb{R}$。因为对于任意 $x$，$\nabla^2 f(x)=P$，所以函数 $f$ 是凸的，当且仅当 $P\succeq0$（$f$ 是凹的当且仅当 $P\preceq0$）。

对于二次函数，严格凸比较容易表达：函数 $f$ 是严格凸的，当且仅当 $P\succ0$（函数是严格凹的当且仅当 $P\prec0$）。

\noindent\textbf{注释 3.1}\quad 在判断函数的凸性和凹性时，不管是一阶条件还是二阶条件，$\dom f$ 必须是凸集这个前提条件必须满足。例如，考虑函数 $f(x)=1/x^2$，其定义域为 $\dom f=\{x\in\mathbb{R}\mid x\neq0\}$，对于所有 $x\in\dom f$ 均满足 $f''(x)>0$，但是函数 $f(x)$ 并不是凸函数。

\subsection*{3.1.5\quad 例子}

前文已经提到所有的线性函数和仿射函数均为凸函数（同时也是凹函数），并描述了凸和凹的一次函数。本节给出更多的凸函数和凹函数的例子。首先考虑 $\mathbb{R}$ 上的一些函数，其自变量为 $x$。

\begin{itemize}
\item \textbf{指数函数。} 对任意 $a\in\mathbb{R}$，函数 $e^{ax}$ 在 $\mathbb{R}$ 上是凸的。
\item \textbf{幂函数。} 当 $a\geq1$ 或 $a\leq0$ 时，$x^a$ 在 $\mathbb{R}_{++}$ 上是凸函数，当 $0\leq a\leq1$ 时，$x^a$ 在 $\mathbb{R}_{++}$ 上是凹函数。
\item \textbf{绝对值幂函数。} 当 $p\geq1$ 时，函数 $|x|^p$ 在 $\mathbb{R}$ 上是凸函数。
\item \textbf{对数函数。} 函数 $\log x$ 在 $\mathbb{R}_{++}$ 上是凹函数。
\item \textbf{负熵。} 函数 $x\log x$ 在其定义域上是凸函数。（定义域为 $\mathbb{R}_{++}$ 或者 $\mathbb{R}_+$，当 $x=0$ 时定义函数值为 0。）
\end{itemize}

我们可以通过基本不等式 (3.1) 或者二阶导数半正定或半负定来判断上述函数是凸的或是凹的。以函数 $f(x)=x\log x$ 为例，其导数和二阶导数为
\[
f'(x)=\log x+1,\qquad f''(x)=1/x,
\]
即对于 $x>0$，有 $f''(x)>0$。所以负熵函数是（严格）凸的。

下面我们给出 $\mathbb{R}^n$ 上的一些例子。

\begin{itemize}
\item \textbf{范数。} $\mathbb{R}^n$ 上的任意范数均为凸函数。
\item \textbf{最大值函数。} 函数 $f(x)=\max\{x_1,\cdots,x_n\}$ 在 $\mathbb{R}^n$ 上是凸的。
\item \textbf{二次-线性分式函数。} 函数 $f(x,y)=x^2/y$，其定义域为
\[
\dom f=\mathbb{R}\times\mathbb{R}_{++}=\{(x,y)\in\mathbb{R}^2\mid y>0\},
\]
是凸函数（如图 3.3 所示）。
\item \textbf{指数和的对数。} 函数 $f(x)=\log(e^{x_1}+\cdots+e^{x_n})$ 在 $\mathbb{R}^n$ 上是凸函数。这个函数可以看成最大值函数的可微（实际上是解析）近似，因为对任意 $x$，下面的不等式成立
\[
\max\{x_1,\cdots,x_n\}\leq f(x)\leq \max\{x_1,\cdots,x_n\}+\log n.
\]
（第二个不等式当 $x$ 的所有分量都相等时是紧的。）图 3.4 描述了当 $n=2$ 时 $f$ 的图像。
\item \textbf{几何平均。} 几何平均函数 $f(x)=\left(\prod_{i=1}^n x_i\right)^{1/n}$ 在定义域 $\dom f=\mathbb{R}_{++}^n$ 上是凹函数。
\item \textbf{对数-行列式。} 函数 $f(X)=\log\det X$ 在定义域 $\dom f=\mathbb{S}_{++}^n$ 上是凹函数。
\end{itemize}

判断上述函数的凸性（或者凹性）可以有多种途径，可以直接验证不等式 (3.1) 是否成立，亦可以验证其 Hessian 矩阵是否半正定，或者可以将函数转换到与其定义域相交的任意直线上，通过得到的单变量函数判断原函数的凸性。

\noindent\textbf{范数。}\quad 如果函数 $f:\mathbb{R}^n\to\mathbb{R}$ 是范数，任取 $0\leq\theta\leq1$，有
\[
f(\theta x+(1-\theta)y)\leq f(\theta x)+f((1-\theta)y)=\theta f(x)+(1-\theta)f(y).
\]
上述不等式可以由三角不等式得到，当范数满足齐次性时，上述不等式取等号。

\noindent\textbf{最大值函数。}\quad 对任意 $0\leq\theta\leq1$，函数 $f(x)=\max_i x_i$ 满足
\[
\begin{aligned}
f(\theta x+(1-\theta)y)
&=\max_i(\theta x_i+(1-\theta)y_i)\\
&\leq \theta\max_i x_i+(1-\theta)\max_i y_i\\
&=\theta f(x)+(1-\theta)f(y).
\end{aligned}
\]

\noindent\textbf{二次-线性分式函数。}\quad 为了说明二次-线性分式函数 $f(x,y)=x^2/y$ 是凸的，我们注意到，对于 $y>0$，有
\[
\nabla^2 f(x,y)=\frac{2}{y^3}
\begin{bmatrix}
y^2 & -xy\\
-xy & x^2
\end{bmatrix}
=\frac{2}{y^3}
\begin{bmatrix}
y\\
-x
\end{bmatrix}
\begin{bmatrix}
y\\
-x
\end{bmatrix}^T\succeq0.
\]

\noindent\textbf{指数和的对数。}\quad 指数和的对数函数的 Hessian 矩阵为
\[
\nabla^2 f(x)=\frac{1}{(\mathbf{1}^Tz)^2}\left((\mathbf{1}^Tz)\diag(z)-zz^T\right),
\]
其中 $z=(e^{x_1},\cdots,e^{x_n})$。为了说明 $\nabla^2 f(x)\succeq0$，我们证明对任意 $v$，有 $v^T\nabla^2 f(x)v\geq0$，即
\[
v^T\nabla^2 f(x)v
=\frac{1}{(\mathbf{1}^Tz)^2}\left(\left(\sum_{i=1}^n z_i\right)\left(\sum_{i=1}^n v_i^2z_i\right)-\left(\sum_{i=1}^n v_iz_i\right)^2\right)\geq0.
\]
上述不等式可以应用 Cauchy-Schwarz 不等式 $(a^Ta)(b^Tb)\geq(a^Tb)^2$ 得到，此时向量 $a$ 和 $b$ 的分量为 $a_i=v_i\sqrt{z_i}$，$b_i=\sqrt{z_i}$。

\noindent\textbf{几何平均。}\quad 类似地，我们说明，几何平均函数 $f(x)=\left(\prod_{i=1}^n x_i\right)^{1/n}$ 在定义域 $\dom f=\mathbb{R}_{++}^n$ 上是凹的。其 Hessian 矩阵 $\nabla^2 f(x)$ 可以通过下面两个式子给出
\[
\frac{\partial^2 f(x)}{\partial x_k^2}=-(n-1)\frac{\left(\prod_{i=1}^n x_i\right)^{1/n}}{n^2x_k^2},\qquad
\frac{\partial^2 f(x)}{\partial x_k\partial x_l}=
\frac{\left(\prod_{i=1}^n x_i\right)^{1/n}}{n^2x_kx_l}\quad \forall k\neq l.
\]
因此 $\nabla^2 f(x)$ 具有如下表达式
\[
\nabla^2 f(x)=-\frac{\prod_{i=1}^n x_i^{1/n}}{n^2}\left(n\diag(1/x_1^2,\cdots,1/x_n^2)-qq^T\right),
\]
其中，$q_i=1/x_i$。我们需要证明 $\nabla^2 f(x)\preceq0$，即对于任意向量 $v$，有
\[
v^T\nabla^2 f(x)v
=-\frac{\prod_{i=1}^n x_i^{1/n}}{n^2}\left(n\sum_{i=1}^n v_i^2/x_i^2-\left(\sum_{i=1}^n v_i/x_i\right)^2\right)\leq0.
\]
这同样可以应用 Cauchy-Schwarz 不等式 $(a^Ta)(b^Tb)\geq(a^Tb)^2$ 得到，只需令向量 $a=\mathbf{1}$，向量 $b$ 的分量 $b_i=v_i/x_i$。

\noindent\textbf{对数-行列式。}\quad 对于函数 $f(X)=\log\det X$，我们可以将其转化为任意直线上的单变量函数来验证它是凹的。令 $X=Z+tV$，其中 $Z,V\in\mathbb{S}^n$，定义 $g(t)=f(Z+tV)$，自变量 $t$ 满足 $Z+tV\succ0$。不失一般性，假设 $t=0$ 满足条件，即 $Z\succ0$。我们有
\[
\begin{aligned}
g(t)&=\log\det(Z+tV)\\
&=\log\det\left(Z^{1/2}(I+tZ^{-1/2}VZ^{-1/2})Z^{1/2}\right)\\
&=\sum_{i=1}^n\log(1+t\lambda_i)+\log\det Z,
\end{aligned}
\]
其中 $\lambda_1,\cdots,\lambda_n$ 是矩阵 $Z^{-1/2}VZ^{-1/2}$ 的特征值。因此下式成立
\[
g'(t)=\sum_{i=1}^n\frac{\lambda_i}{1+t\lambda_i},\qquad
g''(t)=-\sum_{i=1}^n\frac{\lambda_i^2}{(1+t\lambda_i)^2}.
\]
因为 $g''(t)\leq0$，函数 $f$ 是凹的。

\subsection*{3.1.6\quad 下水平集}

函数 $f:\mathbb{R}^n\to\mathbb{R}$ 的 $\alpha$-下水平集定义为
\[
C_\alpha=\{x\in\dom f\mid f(x)\leq\alpha\}.
\]
对于任意 $\alpha$ 值，凸函数的下水平集仍然是凸集。证明可以由凸集的定义直接得到：如果 $x,y\in C_\alpha$，则有 $f(x)\leq\alpha$，$f(y)\leq\alpha$，因此对于任意 $0\leq\theta\leq1$，$f(\theta x+(1-\theta)y)\leq\alpha$，即 $\theta x+(1-\theta)y\in C_\alpha$。

反过来不一定正确：某个函数的所有下水平集都是凸集，但这个函数可能不是凸函数。例如，函数 $f(x)=-e^x$ 在 $\mathbb{R}$ 上不是凸函数（实质上，它是严格凹函数），但是其所有下水平集均为凸集。

如果 $f$ 是凹函数，则由 $\{x\in\dom f\mid f(x)\geq\alpha\}$ 定义的 $\alpha$-上水平集也是凸集。下水平集的性质可以用来判断集合的凸性，若某个集合可以描述为一个凸函数的下水平集，或者一个凹函数的上水平集，则其是凸集。

\noindent\textbf{例 3.3}\quad 对于 $x\in\mathbb{R}_+^n$，其几何平均和算术平均分别为
\[
G(x)=\left(\prod_{i=1}^n x_i\right)^{1/n},\qquad
A(x)=\frac{1}{n}\sum_{i=1}^n x_i,
\]
（在 $G$ 中，我们定义 $0^{1/n}=0$）。算术几何平均不等式为 $G(x)\leq A(x)$。

设 $0\leq\alpha\leq1$，考虑集合
\[
\{x\in\mathbb{R}_+^n\mid G(x)\geq \alpha A(x)\},
\]
即使得几何平均至少大于等于算术平均的 $\alpha$ 倍的集合。此集合是凸集，因为它是凹函数 $G(x)-\alpha A(x)$ 的 $0$-上水平集。事实上，这个集合是正齐次的，因此它是凸锥。

\subsection*{3.1.7\quad 上境图}

函数 $f:\mathbb{R}^n\to\mathbb{R}$ 的图像定义为
\[
\{(x,f(x))\mid x\in\dom f\},
\]
它是 $\mathbb{R}^{n+1}$ 空间的一个子集。函数 $f:\mathbb{R}^n\to\mathbb{R}$ 的上境图定义为
\[
\epi f=\{(x,t)\mid x\in\dom f,\ f(x)\leq t\},
\]
它也是 $\mathbb{R}^{n+1}$ 空间的一个子集。（“Epi” 是之上的意思，所以上境图的英文 epigraph 是“在函数图像之上”的意思。）图 3.5 说明了上境图的定义。

凸集和凸函数的联系可以通过上境图来建立：一个函数是凸函数，当且仅当其上境图是凸集。一个函数是凹函数，当且仅当其亚图
\[
\hypo f=\{(x,t)\mid t\leq f(x)\},
\]
是凸集。

\noindent\textbf{例 3.4\quad 矩阵分式函数。}\quad 矩阵分式函数 $f:\mathbb{R}^n\times\mathbb{S}^n\to\mathbb{R}$，
\[
f(x,Y)=x^TY^{-1}x
\]
在定义域 $\dom f=\mathbb{R}^n\times\mathbb{S}_{++}^n$ 上是凸的。（这实质上是定义域为 $\dom f=\mathbb{R}\times\mathbb{R}_{++}$ 的二次-线性分式函数 $f(x,y)=x^2/y$ 的一个扩展。）

一个简单的方式来验证矩阵分式函数 $f$ 的凸性是通过其上境图
\[
\begin{aligned}
\epi f
&=\{(x,Y,t)\mid Y\succ0,\ x^TY^{-1}x\leq t\}\\
&=\left\{(x,Y,t)\ \middle|\ 
\begin{bmatrix}
Y & x\\
x^T & t
\end{bmatrix}\succeq0,\ Y\succ0
\right\},
\end{aligned}
\]
并应用 Schur 补条件判断分块矩阵的半正定性（见 \S A.5.5）。最后一个条件是关于 $(x,Y,t)$ 的线性矩阵不等式，因此 $\epi f$ 是凸集。

对于 $n=1$ 的特殊情况，矩阵分式函数简化为二次-线性分式函数 $x^2/y$；相应的线性矩阵不等式表示为
\[
\begin{bmatrix}
y & x\\
x & t
\end{bmatrix}\succeq0,\qquad y>0
\]
（函数图像如图 3.3 所示）。

关于凸函数的很多结果可以从几何的角度利用上境图并结合凸集的一些结论来证明（或理解），作为一个例子，考虑凸函数的一阶条件
\[
f(y)\geq f(x)+\nabla f(x)^T(y-x),
\]
其中函数 $f$ 是凸的，$x,y\in\dom f$。我们可以利用 $\epi f$ 从几何角度理解上述基本不等式。如果 $(y,t)\in\epi f$，有
\[
t\geq f(y)\geq f(x)+\nabla f(x)^T(y-x).
\]
上式可以描述为
\[
(y,t)\in\epi f\Longrightarrow
\begin{bmatrix}
\nabla f(x)\\
-1
\end{bmatrix}^T
\left(
\begin{bmatrix}
y\\
t
\end{bmatrix}
-
\begin{bmatrix}
x\\
f(x)
\end{bmatrix}
\right)\leq0.
\]
这意味着法向量为 $(\nabla f(x),-1)$ 的超平面在边界点 $(x,f(x))$ 支撑着 $\epi f$；如图 3.6 所示。

\subsection*{3.1.8\quad Jensen 不等式及其扩展}

基本不等式
\[
f(\theta x+(1-\theta)y)\leq \theta f(x)+(1-\theta)f(y),
\]
有时也称作 Jensen 不等式。此不等式可以很方便地扩展至更多点的凸组合：如果函数 $f$ 是凸函数，$x_1,\cdots,x_k\in\dom f$，$\theta_1,\cdots,\theta_k\geq0$ 且 $\theta_1+\cdots+\theta_k=1$，则下式成立
\[
f(\theta_1x_1+\cdots+\theta_kx_k)\leq \theta_1 f(x_1)+\cdots+\theta_k f(x_k).
\]

考虑凸集时，此不等式可以扩展至无穷项和、积分以及期望。例如，如果在 $S\subseteq\dom f$ 上 $p(x)\geq0$ 且 $\int_S p(x)\,dx=1$，则当相应的积分存在时，下式成立
\[
f\left(\int_S p(x)x\,dx\right)\leq \int_S f(x)p(x)\,dx.
\]
扩展到更一般的情况，我们可以采用其支撑属于 $\dom f$ 的任意概率测度。如果 $x$ 是随机变量，事件 $x\in\dom f$ 发生的概率为 $1$，函数 $f$ 是凸函数，当相应的期望存在时，我们有
\begin{equation}
f(\mathbf{E}x)\leq \mathbf{E}f(x).
\tag{3.5}
\end{equation}
设随机变量 $x$ 的可能取值为 $\{x_1,x_2\}$，相应地取值概率为 $\operatorname{prob}(x=x_1)=\theta$，$\operatorname{prob}(x=x_2)=1-\theta$，则由一般形式 (3.5) 可以得到基本不等式 (3.1)。所以不等式 (3.5) 可以刻画凸性：如果函数 $f$ 不是凸函数，那么存在随机变量 $x$，$x\in\dom f$ 以概率 $1$ 发生，使得 $f(\mathbf{E}x)>\mathbf{E}f(x)$。

上述所有不等式均被称为 Jensen 不等式，而实际上最初由 Jensen 提出的不等式相当简单
\[
f\left(\frac{x+y}{2}\right)\leq \frac{f(x)+f(y)}{2}.
\]

\noindent\textbf{注释 3.2}\quad 我们可以这样理解式 (3.5)。假设 $x\in\dom f\subseteq\mathbb{R}^n$，$z$ 是 $\mathbb{R}^n$ 中的随机变量，其均值为零，则有
\[
\mathbf{E}f(x+z)\geq f(x).
\]
因此，随机化或者扰动（即在自变量上增加一个零均值的随机变量）从平均效果上不会减小凸函数的值。

\subsection*{3.1.9\quad 不等式}

很多著名的不等式都可以通过将 Jensen 不等式应用于合适的凸函数得到。（事实上，凸性和 Jensen 不等式可以构成不等式理论的基础。）作为一个简单的例子，考虑算术-几何平均不等式
\begin{equation}
\sqrt{ab}\leq (a+b)/2,
\tag{3.6}
\end{equation}
其中 $a,b\geq0$。我们可以利用凸性和 Jensen 不等式得到此不等式。函数 $-\log x$ 是凸函数；利用 Jensen 不等式，令 $\theta=1/2$，可得
\[
-\log\left(\frac{a+b}{2}\right)\leq \frac{-\log a-\log b}{2}.
\]
等式两边取指数即可得到式 (3.6)。

作为另一个小例子，我们来证明 Hölder 不等式：对 $p>1$，$1/p+1/q=1$，以及 $x,y\in\mathbb{R}^n$ 有
\[
\sum_{i=1}^n x_i y_i
\leq
\left(\sum_{i=1}^n |x_i|^p\right)^{1/p}
\left(\sum_{i=1}^n |y_i|^q\right)^{1/q}.
\]
由 $-\log x$ 的凸性以及 Jensen 不等式，我们可以得到更为一般的算术-几何平均不等式
\[
a^\theta b^{1-\theta}\leq \theta a+(1-\theta)b,
\]
其中 $a,b\geq0$，$0\leq\theta\leq1$。令
\[
a=\frac{|x_i|^p}{\sum_{j=1}^n |x_j|^p},\qquad
b=\frac{|y_i|^q}{\sum_{j=1}^n |y_j|^q},\qquad
\theta=1/p,
\]
可以得到如下不等式
\[
\left(
\frac{|x_i|^p}{\sum_{j=1}^n |x_j|^p}
\right)^{1/p}
\left(
\frac{|y_i|^q}{\sum_{j=1}^n |y_j|^q}
\right)^{1/q}
\leq
\frac{|x_i|^p}{p\sum_{j=1}^n |x_j|^p}
+
\frac{|y_i|^q}{q\sum_{j=1}^n |y_j|^q}.
\]
对 $i$ 进行求和可以得到 Hölder 不等式。

\section*{3.2\quad 保凸运算}

本节讨论几种保持函数凸性或者凹性的运算，这样可以构造新的凸函数或者凹函数。我们首先从一些简单的运算开始，如求和、伸缩以及逐点上确界，之后再介绍一些更为复杂的运算（其中一些运算的特例即为简单运算）。

\subsection*{3.2.1\quad 非负加权求和}

显而易见，如果函数 $f$ 是凸函数且 $\alpha\geq0$，则函数 $\alpha f$ 也为凸函数。如果函数 $f_1$ 和 $f_2$ 都是凸函数，则它们的和 $f_1+f_2$ 也是凸函数。将非负伸缩以及求和运算结合起来，可以看出，凸函数的集合本身是一个凸锥：凸函数的非负加权求和仍然是凸函数，即函数
\[
f=w_1f_1+\cdots+w_mf_m
\]
是凸函数。类似地，凹函数的非负加权求和仍然是凹函数。严格凸（凹）函数的非负、非零加权求和是严格凸（凹）函数。

这个性质可以扩展至无限项的求和以及积分的情形。例如，如果固定任意 $y\in\mathcal{A}$，函数 $f(x,y)$ 关于 $x$ 是凸函数，且对任意 $y\in\mathcal{A}$，有 $w(y)\geq0$，则函数 $g$
\[
g(x)=\int_{\mathcal{A}} w(y)f(x,y)\,dy
\]
关于 $x$ 是凸函数（若此积分存在）。

我们可以很容易直接验证非负伸缩以及求和运算是保凸运算，或者可以根据相关的上境图得到此结论。例如，如果 $w\geq0$ 且 $f$ 是凸函数，我们有
\[
\epi(wf)=
\begin{bmatrix}
I & 0\\
0 & w
\end{bmatrix}
\epi f,
\]
因为凸集通过线性变换得到的像仍然是凸集，所以 $\epi(wf)$ 是凸集。

\subsection*{3.2.2\quad 复合仿射映射}

假设函数 $f:\mathbb{R}^n\to\mathbb{R}$，$A\in\mathbb{R}^{n\times m}$，以及 $b\in\mathbb{R}^n$，定义 $g:\mathbb{R}^m\to\mathbb{R}$ 为
\[
g(x)=f(Ax+b),
\]
其中 $\dom g=\{x\mid Ax+b\in\dom f\}$。若函数 $f$ 是凸函数，则函数 $g$ 是凸函数；如果函数 $f$ 是凹函数，那么函数 $g$ 是凹函数。

\subsection*{3.2.3\quad 逐点最大和逐点上确界}

如果函数 $f_1$ 和 $f_2$ 均为凸函数，则二者的逐点最大函数 $f$
\[
f(x)=\max\{f_1(x),f_2(x)\},
\]
其定义域为 $\dom f=\dom f_1\cap\dom f_2$，仍然是凸函数。这个性质可以很容易验证：任取 $0\leq\theta\leq1$ 以及 $x,y\in\dom f$，有
\[
\begin{aligned}
f(\theta x+(1-\theta)y)
&=\max\{f_1(\theta x+(1-\theta)y),f_2(\theta x+(1-\theta)y)\}\\
&\leq \max\{\theta f_1(x)+(1-\theta)f_1(y),\theta f_2(x)+(1-\theta)f_2(y)\}\\
&\leq \theta\max\{f_1(x),f_2(x)\}+(1-\theta)\max\{f_1(y),f_2(y)\}\\
&=\theta f(x)+(1-\theta)f(y),
\end{aligned}
\]
从而说明了函数 $f$ 的凸性。同样很容易证明，如果函数 $f_1,\cdots,f_m$ 为凸函数，则它们的逐点最大函数
\[
f(x)=\max\{f_1(x),\cdots,f_m(x)\}
\]
仍然是凸函数。

\noindent\textbf{例 3.5\quad 分片线性函数。}\quad 函数
\[
f(x)=\max\{a_1^Tx+b_1,\cdots,a_L^Tx+b_L\}
\]
定义了一个分片线性（实际上是仿射）函数（具有 $L$ 个或者更少的子区域）。因为它是一系列仿射函数的逐点最大函数，所以它是凸函数。

反之亦成立：任意具有 $L$ 个或者更少子区域的分片线性凸函数都可以表达成上述形式。（见习题 3.29。）

\noindent\textbf{例 3.6\quad 最大 $r$ 个分量之和。}\quad 对于任意 $x\in\mathbb{R}^n$，用 $x_{[i]}$ 表示 $x$ 中第 $i$ 大的分量，即将 $x$ 的分量按照非升序进行排列得到下式
\[
x_{[1]}\geq x_{[2]}\geq \cdots\geq x_{[n]}.
\]
则对 $x$ 的最大 $r$ 个分量进行求和所得到的函数
\[
f(x)=\sum_{i=1}^r x_{[i]},
\]
是凸函数。事实上，此函数可以表达为
\[
f(x)=\sum_{i=1}^r x_{[i]}=
\max\{x_{i_1}+\cdots+x_{i_r}\mid 1\leq i_1<i_2<\cdots<i_r\leq n\},
\]
即从 $x$ 的分量中选取 $r$ 个不同分量进行求和的所有可能组合的最大值。因为函数 $f$ 是 $n!/(r!(n-r)!)$ 个线性函数的逐点最大，所以是凸函数。

作为一个扩展，可以证明当 $w_1\geq w_2\geq\cdots\geq w_r\geq0$ 时，函数 $\sum_{i=1}^r w_i x_{[i]}$ 是凸函数。（见习题 3.19。）

逐点最大的性质可以扩展至无限个凸函数的逐点上确界。如果对于任意 $y\in\mathcal{A}$，函数 $f(x,y)$ 关于 $x$ 都是凸的，则函数 $g$
\begin{equation}
g(x)=\sup_{y\in\mathcal{A}} f(x,y)
\tag{3.7}
\end{equation}
关于 $x$ 亦是凸的。此时，函数 $g$ 的定义域为
\[
\dom g=\{x\mid (x,y)\in\dom f\ \forall y\in\mathcal{A},\ \sup_{y\in\mathcal{A}} f(x,y)<\infty\}.
\]
类似地，一系列凹函数的逐点下确界仍然是凹函数。

从上境图的角度理解，一系列函数的逐点上确界函数对应着这些函数上境图的交集：对于函数 $f$，$g$ 以及式 (3.7) 定义的 $\mathcal{A}$，我们有
\[
\epi g=\bigcap_{y\in\mathcal{A}}\epi f(\cdot,y).
\]
因此，函数 $g$ 的凸性可由一系列凸集的交集仍然是凸集得到。

\noindent\textbf{例 3.7\quad 集合的支撑函数。}\quad 令集合 $C\subseteq\mathbb{R}^n$，且 $C\neq\emptyset$，定义集合 $C$ 的支撑函数 $S_C$ 为
\[
S_C(x)=\sup\{x^Ty\mid y\in C\},
\]
（自然地，函数 $S_C$ 的定义域为 $\dom S_C=\{x\mid \sup_{y\in C}x^Ty<\infty\}$）。对于任意 $y\in C$，$x^Ty$ 是 $x$ 的线性函数，所以 $S_C$ 是一系列线性函数的逐点上确界函数，因此是凸函数。

\noindent\textbf{例 3.8\quad 到集合中最远点的距离。}\quad 令集合 $C\subseteq\mathbb{R}^n$，定义点 $x$ 与集合中最远点的距离（范数）为
\[
f(x)=\sup_{y\in C}\|x-y\|,
\]
此函数是凸函数。为了说明这一点，我们注意到，对于任意 $y$，函数 $\|x-y\|$ 关于 $x$ 是凸函数。因为函数 $f$ 是一族凸函数（对应不同的 $y\in C$）的逐点上确界，所以其是凸函数。

\noindent\textbf{例 3.9\quad 以权为变量的最小二乘费用函数。}\quad 令 $a_1,\cdots,a_n\in\mathbb{R}^m$，在加权最小二乘问题中，我们对所有的 $x\in\mathbb{R}^m$ 极小化目标函数 $\sum_{i=1}^n w_i(a_i^Tx-b_i)^2$。我们称 $w_i$ 为权，并允许负的 $w_i$（则目标函数有可能无下界）。

我们定义（最优）加权最小二乘费用函数为
\[
g(w)=\inf_x\sum_{i=1}^n w_i(a_i^Tx-b_i)^2,
\]
其定义域为
\[
\dom g=\left\{w\ \middle|\ \inf_x\sum_{i=1}^n w_i(a_i^Tx-b_i)^2>-\infty\right\}.
\]
因为函数 $g$ 是一族关于 $w$ 的线性函数的下确界（对应于不同的 $x\in\mathbb{R}^m$），它是 $w$ 的凹函数。

至少在部分定义域上，我们可以得到函数 $g$ 的一个显式表达式。令 $W=\diag(w)$ 是一对角阵，其对角线元素为 $w_1,\cdots,w_n$，令 $A\in\mathbb{R}^{n\times m}$，其行向量为 $a_i^T$，我们有
\[
g(w)=\inf_x(Ax-b)^TW(Ax-b)
=\inf_x(x^TA^TWAx-2b^TWAx+b^TWb).
\]
从上式我们可以看出，若 $A^TWA\not\succeq0$，括号里的二次函数关于 $x$ 无下界，故 $g(w)=-\infty$，即 $w\notin\dom g$。当 $A^TWA\succ0$ 时（即定义了一个严格的线性矩阵不等式），通过解析求解二次函数的极小值，我们可以得到函数 $g$ 的一个简单的表达式
\[
\begin{aligned}
g(w)
&=b^TWb-b^TWA(A^TWA)^{-1}A^TWb\\
&=\sum_{i=1}^n w_i b_i^2-\sum_{i=1}^n w_i^2b_i^2a_i^T
\left(\sum_{j=1}^n w_ja_ja_j^T\right)^{-1}a_i.
\end{aligned}
\]
从上述表达式并不能立即得到函数 $g$ 的凹性（不过可以从矩阵分式函数的凸性来推导函数 $g$ 的凹性：见例 3.4）。

\noindent\textbf{例 3.10\quad 对称矩阵的最大特征值。}\quad 定义函数 $f(X)=\lambda_{\max}(X)$，其定义域为 $\dom f=\mathbb{S}^m$，它是凸函数。为了说明这一点，我们将 $f$ 表述为
\[
f(X)=\sup\{y^TXy\mid \|y\|_2=1\},
\]
即针对不同的 $y\in\mathbb{R}^m$ 关于 $X$ 的一族线性函数（即 $y^TXy$）的逐点上确界。

\noindent\textbf{例 3.11\quad 矩阵范数。}\quad 考虑函数 $f(X)=\|X\|_2$，其定义域为 $\dom f=\mathbb{R}^{p\times q}$，其中 $\|\cdot\|_2$ 表示谱范数或者最大奇异值。函数 $f$ 可以表达为
\[
f(X)=\sup\{u^TXv\mid \|u\|_2=1,\ \|v\|_2=1\},
\]
由于它是 $X$ 的一族线性函数的逐点上确界，所以是凸函数。

作为一个推广，假设 $\|\cdot\|_a$ 和 $\|\cdot\|_b$ 分别是 $\mathbb{R}^p$ 和 $\mathbb{R}^q$ 上的范数，定义矩阵 $X\in\mathbb{R}^{p\times q}$ 的诱导范数为
\[
\|X\|_{a,b}=\sup_{v\neq0}\frac{\|Xv\|_a}{\|v\|_b}.
\]
（当两个范数都取 Euclid 范数时，上述定义范数即为谱范数。）诱导范数可以写成
\[
\begin{aligned}
\|X\|_{a,b}
&=\sup\{\|Xv\|_a\mid \|v\|_b=1\}\\
&=\sup\{u^TXv\mid \|u\|_{a*}=1,\ \|v\|_b=1\},
\end{aligned}
\]
其中 $\|\cdot\|_{a*}$ 是 $\|\cdot\|_a$ 的对偶范数，在此我们利用了
\[
\|z\|_a=\sup\{u^Tz\mid \|u\|_{a*}=1\}.
\]
因此 $\|X\|_{a,b}$ 可以表示成 $X$ 的一系列线性函数的上确界，它是凸函数。

\noindent\textbf{表示成一族仿射函数的逐点上确界}

上述例子描述了一个建立函数凸性的好方法：将其表示为一族仿射函数的逐点上确界。除了一个技术条件，反过来也成立：几乎所有的凸函数都可以表示成一族仿射函数的逐点上确界。例如，如果函数 $f:\mathbb{R}^n\to\mathbb{R}$ 是凸函数，其定义域为 $\dom f=\mathbb{R}^n$，我们有
\[
f(x)=\sup\{g(x)\mid g\text{ 仿射},\ g(z)\leq f(z)\ \forall z\}.
\]
换言之，函数 $f$ 是它所有的仿射全局下估计的逐点上确界。下面我们将证明这个结论，当 $\dom f\neq\mathbb{R}^n$ 的情况留作习题（见习题 3.28）。

设函数 $f$ 是凸函数，定义域为 $\dom f=\mathbb{R}^n$，显然下面的不等式成立
\[
f(x)\geq \sup\{g(x)\mid g\text{ 仿射},\ g(z)\leq f(z)\ \forall z\},
\]
因为函数 $g$ 是函数 $f$ 的任意仿射下估计，我们有 $g(x)\leq f(x)$。为了建立等式，我们说明，对于任意 $x\in\mathbb{R}^n$，存在仿射函数 $g$ 是函数 $f$ 的全局下估计，并且满足 $g(x)=f(x)$。毫无疑问，函数 $f$ 的上境图是凸集，因此我们在点 $(x,f(x))$ 处可以找到此凸集的支撑超平面，即存在 $a\in\mathbb{R}^n$，$b\in\mathbb{R}$ 且 $(a,b)\neq0$，使得对任意 $(z,t)\in\epi f$，有
\[
\begin{bmatrix}
a\\
b
\end{bmatrix}^T
\begin{bmatrix}
x-z\\
f(x)-t
\end{bmatrix}
\leq0.
\]
即对任意 $z\in\dom f=\mathbb{R}^n$ 以及所有 $s\geq0$（$(z,t)\in\epi f$ 等价于存在 $s\geq0$ 使得 $t=f(z)+s$），下式成立
\begin{equation}
a^T(x-z)+b(f(x)-f(z)-s)\leq0.
\tag{3.8}
\end{equation}
为了保证不等式 (3.8) 对所有的 $s\geq0$ 均成立，必须有 $b\geq0$。如果 $b=0$，对所有的 $z\in\mathbb{R}^n$，不等式 (3.8) 可以简化为 $a^T(x-z)\leq0$，这意味着 $a=0$，于是和假设 $(a,b)\neq0$ 矛盾。因此 $b>0$，即支撑超平面不是竖直的。

我们知道 $b>0$，因此，对任意 $z$，令 $s=0$，式 (3.8) 可以重新表述为
\[
g(z)=f(x)+(a/b)^T(x-z)\leq f(z).
\]
由此说明函数 $g$ 是函数 $f$ 的一个仿射下估计，并且满足 $g(x)=f(x)$。

\subsection*{3.2.4\quad 复合}

本节给定函数 $h:\mathbb{R}^k\to\mathbb{R}$ 以及 $g:\mathbb{R}^n\to\mathbb{R}^k$，定义复合函数 $f=h\circ g:\mathbb{R}^n\to\mathbb{R}$ 为
\[
f(x)=h(g(x)),\qquad \dom f=\{x\in\dom g\mid g(x)\in\dom h\}.
\]
我们考虑当函数 $f$ 保凸或者保凹时，函数 $h$ 和 $g$ 必须满足的条件。

\noindent\textbf{标量复合}

首先考虑 $k=1$ 的情况，即 $h:\mathbb{R}\to\mathbb{R}$，$g:\mathbb{R}^n\to\mathbb{R}$。仅考虑 $n=1$ 的情况（事实上，将函数限定在与其定义域相交的任意直线上得到的函数决定了原函数的凸性）。为了找出复合规律，首先假设函数 $h$ 和 $g$ 是二次可微的，且 $\dom g=\dom h=\mathbb{R}$。在上述假设下，函数 $f$ 是凸的等价于 $f''\geq0$（即对所有的 $x\in\mathbb{R}$，$f''(x)\geq0$）。复合函数 $f=h\circ g$ 的二阶导数为
\begin{equation}
f''(x)=h''(g(x))g'(x)^2+h'(g(x))g''(x).
\tag{3.9}
\end{equation}
假设函数 $g$ 是凸函数（$g''\geq0$），函数 $h$ 是凸函数且非减（即 $h''\geq0$ 且 $h'\geq0$），从式 (3.9) 可以得出 $f''\geq0$，即函数 $f$ 是凸函数。类似地，由式 (3.9) 可以得出如下结论
\begin{equation}
\begin{array}{ll}
\text{如果 }h\text{ 是凸函数且非减， }g\text{ 是凸函数，则 }f\text{ 是凸函数，}\\
\text{如果 }h\text{ 是凸函数且非增， }g\text{ 是凹函数，则 }f\text{ 是凸函数，}\\
\text{如果 }h\text{ 是凹函数且非减， }g\text{ 是凹函数，则 }f\text{ 是凹函数，}\\
\text{如果 }h\text{ 是凹函数且非增， }g\text{ 是凸函数，则 }f\text{ 是凹函数。}
\end{array}
\tag{3.10}
\end{equation}
上述结论在函数 $g$ 和 $h$ 二次可微且定义域均为 $\mathbb{R}$ 时成立。事实上，对于更一般的情况，如 $n>1$，不再假设函数 $h$ 和 $g$ 可微或者 $\dom g=\mathbb{R}^n$，$\dom h=\mathbb{R}$，一些相似的复合规则仍然成立
\begin{equation}
\begin{array}{ll}
\text{如果 }h\text{ 是凸函数且 }\tilde h\text{ 非减， }g\text{ 是凸函数，则 }f\text{ 是凸函数，}\\
\text{如果 }h\text{ 是凸函数且 }\tilde h\text{ 非增， }g\text{ 是凹函数，则 }f\text{ 是凸函数，}\\
\text{如果 }h\text{ 是凹函数且 }\tilde h\text{ 非减， }g\text{ 是凹函数，则 }f\text{ 是凹函数，}\\
\text{如果 }h\text{ 是凹函数且 }\tilde h\text{ 非增， }g\text{ 是凸函数，则 }f\text{ 是凹函数。}
\end{array}
\tag{3.11}
\end{equation}
其中，$\tilde h$ 表示函数 $h$ 的扩展值延伸，若点不在 $\dom h$ 内，对其赋值 $\infty$（若 $h$ 是凸函数）或者 $-\infty$（若 $h$ 是凹函数）。这些结论和式 (3.10) 中的结论的唯一不同是我们要求扩展值延伸 $\tilde h$ 在整个 $\mathbb{R}$ 上非增或者非减。

为了更好地理解式 (3.11) 中的结论，假设 $h$ 是凸函数，所以 $\tilde h$ 在定义域 $\dom h$ 外取值为 $\infty$。$\tilde h$ 非减意味着对于任意 $x,y\in\mathbb{R}$，$x<y$，我们有 $\tilde h(x)\leq\tilde h(y)$。特别地，若 $y\in\dom h$，则 $x\in\dom h$。换言之，我们可以认为 $h$ 的定义域在负方向上无限延伸；它或者是 $\mathbb{R}$ 或者是形如 $(-\infty,a)$ 或 $(-\infty,a]$ 的区间。类似地，若 $h$ 是凸函数且 $\tilde h$ 非增，我们可以理解为 $h$ 是非增的且 $\dom h$ 在正方向上趋于无穷。图 3.7 描述了不同的扩展值延伸的情况。

\noindent\textbf{例 3.12}\quad 通过一些简单例子，我们可以更好地理解复合定理中函数 $h$ 需要满足的条件。
\begin{itemize}
\item 函数 $h(x)=\log x$，定义域为 $\dom h=\mathbb{R}_{++}$，其为凹函数且 $\tilde h$ 非减。
\item 函数 $h(x)=x^{1/2}$，定义域为 $\dom h=\mathbb{R}_+$，其为凹函数且 $\tilde h$ 非减。
\item 函数 $h(x)=x^{3/2}$，定义域为 $\dom h=\mathbb{R}_+$，其为凸函数，但是不满足 $\tilde h$ 非减的条件。例如，$\tilde h(-1)=\infty$ 但 $\tilde h(1)=1$。
\item 当 $x\geq0$ 时，$h(x)=x^{3/2}$，当 $x<0$ 时 $h(x)=0$，定义域为 $\dom h=\mathbb{R}$，$h$ 是凸函数且满足 $\tilde h$ 非减的条件。
\end{itemize}

即使不假设可微并运用表达式 (3.9)，也可以直接证明复合函数结论式 (3.11)。作为一个例子，我们证明如下结论：如果 $g$ 是凸函数，$h$ 是凸函数且 $h$ 非减，则 $f=h\circ g$ 是凸函数。假设 $x,y\in\dom f$，$0\leq\theta\leq1$。由于 $x,y\in\dom f$，我们有 $x,y\in\dom g$ 且 $g(x),g(y)\in\dom h$。因为 $\dom g$ 是凸集，有 $\theta x+(1-\theta)y\in\dom g$；由函数 $g$ 的凸性可得
\begin{equation}
g(\theta x+(1-\theta)y)\leq \theta g(x)+(1-\theta)g(y).
\tag{3.12}
\end{equation}
由 $g(x),g(y)\in\dom h$ 可得 $\theta g(x)+(1-\theta)g(y)\in\dom h$，即式 (3.12) 的右端在 $\dom h$ 内。根据假设 $\tilde h$ 是非减的，可以理解为其定义域在负方向上无限延伸。由式 (3.12) 的右端在 $\dom h$ 内，我们知道其左侧仍在定义域内，即 $g(\theta x+(1-\theta)y)\in\dom h$，因此 $\dom f$ 是凸集。

根据前提条件，$\tilde h$ 非减，利用不等式 (3.12)，我们有
\begin{equation}
h(g(\theta x+(1-\theta)y))\leq h(\theta g(x)+(1-\theta)g(y)).
\tag{3.13}
\end{equation}
由函数 $h$ 的凸性可得
\begin{equation}
h(\theta g(x)+(1-\theta)g(y))\leq \theta h(g(x))+(1-\theta)h(g(y)).
\tag{3.14}
\end{equation}
综合式 (3.13) 和式 (3.14) 可得
\[
h(g(\theta x+(1-\theta)y))\leq \theta h(g(x))+(1-\theta)h(g(y)),
\]
复合定理得证。

\noindent\textbf{例 3.13\quad 简单的复合结论。}
\begin{itemize}
\item 如果 $g$ 是凸函数则 $\exp g(x)$ 是凸函数。
\item 如果 $g$ 是凹函数且大于零，则 $\log g(x)$ 是凹函数。
\item 如果 $g$ 是凹函数且大于零，则 $1/g(x)$ 是凸函数。
\item 如果 $g$ 是凸函数且不小于零，$p\geq1$，则 $g(x)^p$ 是凸函数。
\item 如果 $g$ 是凸函数，则 $-\log(-g(x))$ 在 $\{x\mid g(x)<0\}$ 上是凸函数。
\end{itemize}

\noindent\textbf{注释 3.3}\quad 扩展值延伸 $\tilde h$ 的单调性要求必须满足，注意到是 $\tilde h$，而不仅仅是 $h$。例如，考虑 $g(x)=x^2$，$\dom g=\mathbb{R}$，$h(x)=0$，$\dom h=[1,2]$ 复合的情形。此时 $g$ 是凸函数，$h$ 是凸函数且非减，但是函数 $f=h\circ g$
\[
f(x)=0,\qquad \dom f=[-\sqrt{2},-1]\cup[1,\sqrt{2}],
\]
不是凸函数，因为其定义域非凸。当然，此时函数 $\tilde h$ 不是非减的。

\noindent\textbf{矢量复合}

下面考虑 $k>1$ 的情况，此时更复杂一些。设
\[
f(x)=h(g(x))=h(g_1(x),\cdots,g_k(x)),
\]
其中 $h:\mathbb{R}^k\to\mathbb{R}$，$g_i:\mathbb{R}^n\to\mathbb{R}$。和上一节一样，不失一般性，我们假设 $n=1$。和 $k=1$ 的情形类似，为了得到复合规则，我们假设函数二次可微，且 $\dom g=\mathbb{R}$，$\dom h=\mathbb{R}^k$。此时，我们对函数 $f$ 进行二次微分可得
\begin{equation}
f''(x)=g'(x)^T\nabla^2 h(g(x))g'(x)+\nabla h(g(x))^Tg''(x),
\tag{3.15}
\end{equation}
上式可以看成式 (3.9) 对应的向量形式。同样我们需要判断在什么条件下对所有 $x$ 有 $f''(x)\geq0$（或者对所有 $x$ 有 $f''(x)\leq0$，此时 $f$ 是凹函数）。利用式 (3.15)，我们可以得到很多规则，例如：
\[
\begin{array}{ll}
\text{如果 }h\text{ 是凸函数且在每维分量上 }h\text{ 非减， }g_i\text{ 是凸函数，则 }f\text{ 是凸函数；}\\
\text{如果 }h\text{ 是凸函数且在每维分量上 }h\text{ 非增， }g_i\text{ 是凹函数，则 }f\text{ 是凸函数；}\\
\text{如果 }h\text{ 是凹函数且在每维分量上 }h\text{ 非减， }g_i\text{ 是凹函数，则 }f\text{ 是凹函数。}
\end{array}
\]
和标量的情形类似，对于更一般的情况：$n>1$，不假设 $h$ 或 $g$ 可微以及一般的定义域，类似的复合结论仍然成立。对于一般的结论，不仅 $h$ 需要满足单调性条件，其扩展值延伸 $\tilde h$ 同样必须满足。

为了更好地理解扩展值延伸 $\tilde h$ 必须满足单调性条件的含义，我们考虑凸函数 $h:\mathbb{R}^k\to\mathbb{R}$，且 $\tilde h$ 非减，即对任意 $u\preceq v$，有 $\tilde h(u)\leq \tilde h(v)$。这说明了如果 $v\in\dom h$，则 $u\in\dom h$；$h$ 的定义域在方向 $-\mathbb{R}_+^k$ 上必须无限延伸。这个条件可以紧凑地描述为
\[
\dom h-\mathbb{R}_+^k=\dom h.
\]

\noindent\textbf{例 3.14\quad 矢量复合的例子。}
\begin{itemize}
\item 令 $h(z)=z_{[1]}+\cdots+z_{[r]}$，即对 $z\in\mathbb{R}^k$ 的前 $r$ 大分量进行求和。则 $h$ 是凸函数且在每一维分量上非减。假设 $g_1,\cdots,g_k$ 是 $\mathbb{R}^n$ 上的凸函数，则复合函数 $f=h\circ g$，即最大 $r$ 个 $g_i$ 函数的逐点和，是凸函数。
\item 函数 $h(z)=\log\left(\sum_{i=1}^k e^{z_i}\right)$ 是凸函数且在每一维分量上非减，因此只要 $g_i$ 是凸函数，$\log\left(\sum_{i=1}^k e^{g_i}\right)$ 就是凸函数。
\item 对 $0<p\leq1$，定义在 $\mathbb{R}_+^k$ 上的函数 $h(z)=\left(\sum_{i=1}^k z_i^p\right)^{1/p}$ 是凹的，且其扩展值延伸（当 $z\nsucceq0$ 时为 $-\infty$）在每维分量上非减，则若 $g_i$ 是凹函数且非负，$f(x)=\left(\sum_{i=1}^k g_i(x)^p\right)^{1/p}$ 是凹函数。
\item 设 $p\geq1$，$g_1,\cdots,g_k$ 是凸函数且非负。则函数 $\left(\sum_{i=1}^k g_i(x)^p\right)^{1/p}$ 是凸函数。为了说明这一点，考虑函数 $h:\mathbb{R}^k\to\mathbb{R}$
\[
h(z)=\left(\sum_{i=1}^k \max\{z_i,0\}^p\right)^{1/p},
\]
其中 $\dom h=\mathbb{R}^k$，因此 $h=\tilde h$。由函数 $h$ 是凸函数且非减可知 $h(g(x))$ 关于 $x$ 是凸函数。对 $z\succeq0$，我们有 $h(z)=\left(\sum_{i=1}^k z_i^p\right)^{1/p}$，所以 $\left(\sum_{i=1}^k g_i(x)^p\right)^{1/p}$ 是凸函数。
\item 几何平均函数 $h(z)=\left(\prod_{i=1}^k z_i\right)^{1/k}$，定义域为 $\mathbb{R}_+^k$，它是凹函数，且其扩展值延伸在每维分量上非减。因此若 $g_1,\cdots,g_k$ 是非负凹函数，它们的几何平均 $\left(\prod_{i=1}^k g_i\right)^{1/k}$ 也是非负凹函数。
\end{itemize}

\subsection*{3.2.5\quad 最小化}

我们已经得到，任意个凸函数的逐点最大或者上确界仍然是凸函数。事实上，一些特殊形式的最小化同样可以得到凸函数。如果函数 $f$ 关于 $(x,y)$ 是凸函数，集合 $C$ 是非空凸集，定义函数
\begin{equation}
g(x)=\inf_{y\in C} f(x,y),
\tag{3.16}
\end{equation}
若存在某个 $x$ 使得 $g(x)>-\infty$（该条件隐含着对所有 $x$，$g(x)>-\infty$），则函数 $g$ 关于 $x$ 是凸函数。函数 $g$ 的定义域是 $\dom f$ 在 $x$ 方向上的投影，即
\[
\dom g=\{x\mid \text{对某个 }y\in C\text{ 成立 }(x,y)\in\dom f\}.
\]
任取 $x_1,x_2\in\dom g$，我们利用 Jensen 不等式来证明上述结论。令 $\epsilon>0$，则存在 $y_1,y_2\in C$，使 $f(x_i,y_i)\leq g(x_i)+\epsilon$（$i=1,2$）。设 $\theta\in[0,1]$，我们有
\[
\begin{aligned}
g(\theta x_1+(1-\theta)x_2)
&=\inf_{y\in C} f(\theta x_1+(1-\theta)x_2,y)\\
&\leq f(\theta x_1+(1-\theta)x_2,\theta y_1+(1-\theta)y_2)\\
&\leq \theta f(x_1,y_1)+(1-\theta)f(x_2,y_2)\\
&\leq \theta g(x_1)+(1-\theta)g(x_2)+\epsilon.
\end{aligned}
\]
因为上式对任意 $\epsilon>0$ 均成立，所以下式成立
\[
g(\theta x_1+(1-\theta)x_2)\leq \theta g(x_1)+(1-\theta)g(x_2).
\]
此结论亦可通过上境图来说明。对式 (3.16) 中定义的 $f$，$g$ 和 $C$，设对每个 $x$，在集合 $y\in C$ 上求下确界均可达到，则有
\[
\epi g=\{(x,t)\mid \text{对某个 }y\in C\text{ 成立 }(x,y,t)\in\epi f\}.
\]
由于 $\epi g$ 是凸集在其中一些分量上的投影，所以它仍然是凸集。

\noindent\textbf{例 3.15\quad Schur 补。}\quad 设二次函数
\[
f(x,y)=x^TAx+2x^TBy+y^TCy,
\]
（其中 $A$ 和 $C$ 是对称矩阵）关于 $(x,y)$ 是凸函数，即
\[
\begin{bmatrix}
A & B\\
B^T & C
\end{bmatrix}\succeq0.
\]
我们可以将 $g(x)=\inf_y f(x,y)$ 表述为
\[
g(x)=x^T(A-BC^\dagger B^T)x,
\]
其中 $C^\dagger$ 是矩阵 $C$ 的伪逆（见 \S A.5.4）。根据极小化的性质，$g$ 是凸函数，因此 $A-BC^\dagger B^T\succeq0$。

如果矩阵 $C$ 可逆，即 $C\succ0$，则矩阵 $A-BC^{-1}B^T$ 称为 $C$ 在矩阵
\[
\begin{bmatrix}
A & B\\
B^T & C
\end{bmatrix}
\]
中的 Schur 补（见 \S A.5.5）。

\noindent\textbf{例 3.16\quad 到某一集合的距离。}\quad 采用范数 $\|\cdot\|$，某点 $x$ 到集合 $S\subseteq\mathbb{R}^n$ 的距离定义为
\[
\dist(x,S)=\inf_{y\in S}\|x-y\|.
\]
函数 $\|x-y\|$ 关于 $(x,y)$ 是凸的，所以若集合 $S$ 是凸集，距离函数 $\dist(x,S)$ 是 $x$ 的凸函数。

\noindent\textbf{例 3.17}\quad 设 $h$ 是凸函数。则函数 $g$
\[
g(x)=\inf\{h(y)\mid Ay=x\}
\]
是凸函数。为了说明这一点，定义函数 $f$
\[
f(x,y)=
\begin{cases}
h(y), & \text{如果 }Ay=x,\\
\infty, & \text{其他情况},
\end{cases}
\]
此函数关于 $(x,y)$ 是凸的。以 $y$ 为自变量，极小化函数 $f$ 即可得到函数 $g$，因此函数 $g$ 是凸函数。（直接证明 $g$ 是凸函数亦不复杂。）

\subsection*{3.2.6\quad 透视函数}

给定函数 $f:\mathbb{R}^n\to\mathbb{R}$，则 $f$ 的透视函数 $g:\mathbb{R}^{n+1}\to\mathbb{R}$ 定义为
\[
g(x,t)=t f(x/t),
\]
其定义域为
\[
\dom g=\{(x,t)\mid x/t\in\dom f,\ t>0\}.
\]
透视运算是保凸运算：如果函数 $f$ 是凸函数，则其透视函数 $g$ 也是凸函数。类似地，若 $f$ 是凹函数，则 $g$ 亦是凹函数。

可以从多个角度来证明此结论，例如，我们可以直接验证定义凸性的不等式（见习题 3.33）。这里应用上境图和 \S2.3.3 所描述的 $\mathbb{R}^{n+1}$ 上的透视映射给出一个简短的证明（同时也可以说明“透视”一词的由来）。当 $t>0$ 时，我们有
\[
(x,t,s)\in\epi g
\Longleftrightarrow t f(x/t)\leq s
\Longleftrightarrow f(x/t)\leq s/t
\Longleftrightarrow (x/t,s/t)\in\epi f.
\]
因此，$\epi g$ 是透视映射下 $\epi f$ 的原像，此透视映射将 $(u,v,w)$ 映射为 $(u,w)/v$。根据 \S2.3.3 中的结论，$\epi g$ 是凸集，所以函数 $g$ 是凸函数。

\noindent\textbf{例 3.18\quad Euclid 范数的平方。}\quad $\mathbb{R}^n$ 上的凸函数 $f(x)=x^Tx$ 的透视函数由下式给出
\[
g(x,t)=t(x/t)^T(x/t)=\frac{x^Tx}{t},
\]
当 $t>0$ 时它关于 $(x,t)$ 是凸函数。

我们可以利用其他方法导出 $g$ 的凸性。首先，将 $g$ 表示为一系列二次-线性分式函数 $x_i^2/t$ 的和。在 \S3.1.5 中，我们已经知道，每一项 $x_i^2/t$ 是凸函数，因此和亦为凸函数。另一方面，我们可以将 $g$ 表述为一种特殊的矩阵分式函数 $x^T(tI)^{-1}x$，由此导出凸性（见例 3.4）。

\noindent\textbf{例 3.19\quad 负对数。}\quad 考虑 $\mathbb{R}_{++}$ 上的凸函数 $f(x)=-\log x$，其透视函数为
\[
g(x,t)=-t\log(x/t)=t\log(t/x)=t\log t-t\log x,
\]
在 $\mathbb{R}_{++}^2$ 上它是凸函数。函数 $g$ 称为关于 $t$ 和 $x$ 的相对熵。当 $x=1$ 时，$g$ 即为负熵函数。

基于函数 $g$ 的凸性，我们可以得出一些有趣的相关函数的凸性或凹性。首先，定义两个向量 $u,v\in\mathbb{R}_{++}^n$ 的相对熵
\[
\sum_{i=1}^n u_i\log(u_i/v_i),
\]
由于它是一系列 $u_i$，$v_i$ 的相对熵的和，因此关于 $(u,v)$ 是凸函数。

另一个密切相关的函数是向量 $u,v\in\mathbb{R}_{++}^n$ 之间的 Kullback-Leibler 散度，其形式为
\begin{equation}
D_{\mathrm{kl}}(u,v)=\sum_{i=1}^n\left(u_i\log(u_i/v_i)-u_i+v_i\right),
\tag{3.17}
\end{equation}
因为它是 $(u,v)$ 的相对熵和线性函数的和，所以它也是凸函数。Kullback-Leibler 散度总是满足 $D_{\mathrm{kl}}(u,v)\geq0$，当且仅当 $u=v$ 时，$D_{\mathrm{kl}}(u,v)=0$，因此 Kullback-Leibler 散度可以用来衡量两个正向量之间的偏差；见习题 3.13。（注意到当 $u$ 和 $v$ 都是概率向量，即 $\mathbf{1}^Tu=\mathbf{1}^Tv=1$ 时，相对熵和 Kullback-Leibler 散度是等价的。）

如果在相对熵函数中选择 $v_i=\mathbf{1}^Tu$，我们可以得到定义在 $u\in\mathbb{R}_{++}^n$ 上的凹函数（也是齐次函数）
\[
\sum_{i=1}^n u_i\log(\mathbf{1}^Tu/u_i)
=(\mathbf{1}^Tu)\sum_{i=1}^n z_i\log(1/z_i),
\]
其中 $z=u/(\mathbf{1}^Tu)$。此函数称为归一化熵函数。向量 $z=u/\mathbf{1}^Tu$ 的分量和为 $1$，称为归一化向量或者概率分布；$u$ 的归一化熵是 $\mathbf{1}^Tu$ 和归一化概率分布 $z$ 的熵的乘积。

\noindent\textbf{例 3.20}\quad 设 $f:\mathbb{R}^m\to\mathbb{R}$ 是凸函数，$A\in\mathbb{R}^{m\times n}$，$b\in\mathbb{R}^m$，$c\in\mathbb{R}^n$，$d\in\mathbb{R}$。定义
\[
g(x)=(c^Tx+d)f((Ax+b)/(c^Tx+d)),
\]
其定义域为
\[
\dom g=\{x\mid c^Tx+d>0,\ (Ax+b)/(c^Tx+d)\in\dom f\}.
\]
则 $g$ 是凸函数。

\section*{3.3\quad 共轭函数}

本节介绍一个运算，它将在后续章节发挥重要的作用。

\subsection*{3.3.1\quad 定义及例子}

设函数 $f:\mathbb{R}^n\to\mathbb{R}$，定义函数 $f^*:\mathbb{R}^n\to\mathbb{R}$ 为
\begin{equation}
f^*(y)=\sup_{x\in\dom f}\left(y^Tx-f(x)\right),
\tag{3.18}
\end{equation}
此函数称为函数 $f$ 的共轭函数。使上述上确界有限，即差值 $y^Tx-f(x)$ 在 $\dom f$ 有上界的所有 $y\in\mathbb{R}^n$ 构成了共轭函数的定义域。图 3.8 描述了此定义。

显而易见，$f^*$ 是凸函数，这是因为它是一系列 $y$ 的凸函数（实质上是仿射函数）的逐点上确界。无论 $f$ 是否是凸函数，$f^*$ 都是凸函数。（注意到这里当 $f$ 是凸函数时，下标 $x\in\dom f$ 可以去掉，这是因为根据之前关于扩展值延伸的定义，对于 $x\notin\dom f$，$y^Tx-f(x)=-\infty$。）

我们从一些简单的例子开始描述共轭函数的一些规律。在此基础上我们可以写出很多常见凸函数的共轭函数的解析形式。

\noindent\textbf{例 3.21}\quad 考虑 $\mathbb{R}$ 上一些凸函数的共轭函数。
\begin{itemize}
\item 仿射函数。设 $f(x)=ax+b$。如果 $y=a$，则 $yx-f(x)=-b$；否则 $yx-f(x)$ 对 $x$ 无上界。因此 $\dom f^*=\{a\}$，且 $f^*(a)=-b$。
\item 负对数函数。设 $f(x)=-\log x$，$\dom f=\mathbb{R}_{++}$。当 $y\geq0$ 时，$xy+\log x$ 对 $x>0$ 无上界。当 $y<0$ 时，$xy+\log x$ 在 $x=-1/y$ 处取最大值。因此 $\dom f^*=-\mathbb{R}_{++}$，且
\[
f^*(y)=-\log(-y)-1,\quad y<0.
\]
\item 指数函数。设 $f(x)=e^x$。当 $y<0$ 时，$xy-e^x$ 对 $x$ 无上界。当 $y\geq0$ 时，最大值在 $x=\log y$ 处取得（对于 $y=0$，最大值是 $0$，但不在有限的 $x$ 处取得）。因此 $\dom f^*=\mathbb{R}_+$，且
\[
f^*(y)=y\log y-y,
\]
其中我们约定 $0\log0=0$。
\item 负熵函数。设 $f(x)=x\log x$，$\dom f=\mathbb{R}_+$（同上面讨论，$f(0)=0$）。函数 $xy-x\log x$ 在 $x=e^{y-1}$ 处取得最大值，因此
\[
f^*(y)=e^{y-1}.
\]
\item 反函数。设 $f(x)=1/x$，$x\in\mathbb{R}_{++}$。当 $y>0$ 时，$xy-1/x$ 无上界；当 $y=0$ 时，其上确界为 $0$，但不在有限的 $x$ 处取得；当 $y<0$ 时，最大值在 $x=1/\sqrt{-y}$ 处取得。因此
\[
f^*(y)=-2(-y)^{1/2},
\]
其定义域为 $-\mathbb{R}_+$。
\end{itemize}

\noindent\textbf{例 3.22\quad 严格凸的二次函数。}\quad 考虑函数 $f(x)=(1/2)x^TQx$，其中 $Q\in\mathbb{S}_{++}^n$。对所有的 $y$，$x$ 的函数
\[
y^Tx-\frac{1}{2}x^TQx
\]
都有上界并在 $x=Q^{-1}y$ 处达到上确界，因此
\[
f^*(y)=\frac{1}{2}y^TQ^{-1}y.
\]

\noindent\textbf{例 3.23\quad 对数-行列式。}\quad 我们考虑 $\mathbb{S}_{++}^n$ 上定义的函数 $f(X)=\log\det X^{-1}$。其共轭函数定义为
\[
f^*(Y)=\sup_{X\succ0}\left(\operatorname{tr}(YX)+\log\det X\right),
\]
其中，$\operatorname{tr}(YX)$ 是 $\mathbb{S}^n$ 上的标准内积。首先我们说明只有当 $Y\prec0$ 时，$\operatorname{tr}(YX)+\log\det X$ 才有上界。如果 $Y\nprec0$，则 $Y$ 有特征向量 $v$，$\|v\|_2=1$ 且对应的特征值 $\lambda\geq0$。令 $X=I+tvv^T$，我们有
\[
\operatorname{tr}(YX)+\log\det X
=\operatorname{tr}Y+t\lambda+\log\det(I+tvv^T)
=\operatorname{tr}Y+t\lambda+\log(1+t),
\]
当 $t\to\infty$ 时，上式无界。

接下来考虑 $Y\prec0$ 的情形。为了求最大值，令对 $X$ 的偏导为零，则
\[
\nabla_X\left(\operatorname{tr}(YX)+\log\det X\right)=Y+X^{-1}=0
\]
（见 \S A.4.1），得 $X=-Y^{-1}$（$X$ 是正定的）。因此
\[
f^*(Y)=\log\det(-Y)^{-1}-n,
\]
其定义域为 $\dom f^*=-\mathbb{S}_{++}^n$。

\noindent\textbf{例 3.24\quad 示性函数。}\quad 设 $I_S$ 是某个集合 $S\subseteq\mathbb{R}^n$（不一定是凸集）的示性函数，即当 $x$ 在 $\dom I_S=S$ 内时，$I_S(x)=0$。示性函数的共轭函数为
\[
I_S^*(y)=\sup_{x\in S}y^Tx,
\]
它是集合 $S$ 的支撑函数。

\noindent\textbf{例 3.25\quad 指数和的对数函数。}\quad 为了得到指数和的对数函数
\[
f(x)=\log\left(\sum_i e^{x_i}\right)
\]
的共轭函数，首先考察 $y$ 取何值时 $y^Tx-f(x)$ 的最大值可以得到。对 $x$ 求导，令其为零，我们得到如下条件
\[
y_i=\frac{e^{x_i}}{\sum_j e^{x_j}},\quad i=1,\ldots,n.
\]
当且仅当 $y\succeq0$ 以及 $\mathbf{1}^Ty=1$ 时上述方程有解。将 $y_i$ 的表达式代入 $y^Tx-f(x)$ 可得
\[
f^*(y)=\sum_i y_i\log y_i.
\]
根据前面的约定，$0\log0=0$，故这一表达式在边界点也成立。

反之，如果 $y$ 有负分量，比如 $y_i<0$，令 $x_i\to-\infty$，并保持其他分量不变，则 $y^Tx-f(x)$ 无界。如果 $\mathbf{1}^Ty\neq1$，则令 $x=t\mathbf{1}$，得到
\[
y^Tx-f(x)=t\mathbf{1}^Ty-\log\left(e^t\sum_i1\right)
=t(\mathbf{1}^Ty-1)-\log n,
\]
当 $t\to\infty$ 或 $t\to-\infty$ 时无界。总之，
\[
f^*(y)=
\begin{cases}
\sum_i y_i\log y_i, & y\succeq0,\ \mathbf{1}^Ty=1,\\
\infty, & \text{其他情况}.
\end{cases}
\]
换言之，指数和的对数函数的共轭函数是概率单纯形内的负熵函数。

\noindent\textbf{例 3.26\quad 范数。}\quad 令 $\|\cdot\|$ 表示 $\mathbb{R}^n$ 上的范数，其对偶范数为 $\|\cdot\|_*$。我们说明 $f(x)=\|x\|$ 的共轭函数为
\[
f^*(y)=
\begin{cases}
0, & \|y\|_*\leq1,\\
\infty, & \text{其他情况},
\end{cases}
\]
即范数的共轭函数是对偶范数单位球的示性函数。

如果 $\|y\|_*>1$，根据对偶范数定义，存在 $z\in\mathbb{R}^n$，$\|z\|\leq1$ 使 $y^Tz>1$。取 $x=tz$，$t\to\infty$，得
\[
y^Tx-\|x\|=t(y^Tz-\|z\|)\to\infty,
\]
即 $f^*(y)=\infty$。反之，若 $\|y\|_*\leq1$，对任意 $x$，
\[
y^Tx\leq\|x\|\|y\|_*\leq\|x\|,
\]
所以 $y^Tx-\|x\|\leq0$，且在 $x=0$ 处达到最大值 $0$。

\noindent\textbf{例 3.27\quad 范数的平方。}\quad 考虑
\[
f(x)=\frac{1}{2}\|x\|^2.
\]
我们说明其共轭函数为
\[
f^*(y)=\frac{1}{2}\|y\|_*^2.
\]
根据对偶范数的定义，
\[
y^Tx-\frac{1}{2}\|x\|^2\leq\|y\|_*\|x\|-\frac{1}{2}\|x\|^2.
\]
上式右边关于 $\|x\|$ 的最大值为 $(1/2)\|y\|_*^2$。另一方面，存在一列 $z_k$ 满足 $\|z_k\|\leq1$ 且 $y^Tz_k\to\|y\|_*$，令 $x=\|y\|_*z_k$，可知上界可以任意逼近。因此 $f^*(y)=(1/2)\|y\|_*^2$。

\noindent\textbf{例 3.28\quad 总收入和收益函数。}\quad 考虑某个公司，它消耗资源 $r$，并产生销售总收入 $S(r)$。若资源价格为 $p$，购买资源的成本为 $p^Tr$，于是利润为 $S(r)-p^Tr$。最大利润函数为
\[
M(p)=\sup_r\left(S(r)-p^Tr\right).
\]
显然，
\[
M(p)=(-S)^*(-p).
\]
因此最大收益（资源价格函数）和销售总收入（资源消耗量函数）的共轭密切相关。

\subsection*{3.3.2\quad 基本性质}

\noindent\textbf{Fenchel 不等式}\quad 从共轭函数的定义我们可以得到，对任意 $x$ 和 $y$，如下不等式成立
\[
f(x)+f^*(y)\geq x^Ty,
\]
上述不等式即为 Fenchel 不等式（当 $f$ 可微的时候亦称为 Young 不等式）。以函数 $f(x)=(1/2)x^TQx$ 为例，其中 $Q\in\mathbb{S}_{++}^n$，我们可以得到如下不等式
\[
x^Ty\leq(1/2)x^TQx+(1/2)y^TQ^{-1}y.
\]

\noindent\textbf{共轭的共轭}\quad 上面的例子以及“共轭”的名称都隐含了凸函数的共轭函数的共轭函数是原函数。也即：如果函数 $f$ 是凸函数且 $f$ 是闭的（即 $\epi f$ 是闭集，见 \S A.3.3），则 $f^{**}=f$。例如，若 $\dom f=\mathbb{R}^n$，则我们有 $f^{**}=f$，即 $f$ 的共轭函数的共轭函数还是 $f$（见习题 3.39）。

\noindent\textbf{可微函数}\quad 可微函数 $f$ 的共轭函数亦称为函数 $f$ 的 Legendre 变换。（为了区分一般情况和可微情况下所定义的共轭，一般函数的共轭有时称为 Fenchel 共轭。）

设函数 $f$ 是凸函数且可微，其定义域为 $\dom f=\mathbb{R}^n$，使 $y^Tx-f(x)$ 取最大的 $x^*$ 满足 $y=\nabla f(x^*)$；反之，若 $x^*$ 满足 $y=\nabla f(x^*)$，$y^Tx-f(x)$ 在 $x^*$ 处取最大值。因此，如果 $y=\nabla f(x^*)$，我们有
\[
f^*(y)=x^{*T}\nabla f(x^*)-f(x^*).
\]
所以，给定任意 $y$，我们可以求解梯度方程 $y=\nabla f(x)$，从而得到 $y$ 处的共轭函数 $f^*(y)$。

我们亦可以换一个角度理解。任选 $z\in\mathbb{R}^n$，令 $y=\nabla f(z)$，则
\[
f^*(y)=z^T\nabla f(z)-f(z).
\]

\noindent\textbf{伸缩变换和复合仿射变换}\quad 若 $a>0$ 以及 $b\in\mathbb{R}$，$g(x)=af(x)+b$ 的共轭函数为
\[
g^*(y)=af^*(y/a)-b.
\]
设 $A\in\mathbb{R}^{n\times n}$ 非奇异，$b\in\mathbb{R}^n$，则函数 $g(x)=f(Ax+b)$ 的共轭函数为
\[
g^*(y)=f^*(A^{-T}y)-b^TA^{-T}y,
\]
其定义域为 $\dom g^*=A^T\dom f^*$。

\noindent\textbf{独立函数的和}\quad 如果函数 $f(u,v)=f_1(u)+f_2(v)$，其中 $f_1$ 和 $f_2$ 是凸函数，且共轭函数分别为 $f_1^*$ 和 $f_2^*$，则
\[
f^*(w,z)=f_1^*(w)+f_2^*(z).
\]
换言之，独立凸函数的和的共轭函数是各个凸函数的共轭函数的和。（“独立”的含义是各个函数具有不同的变量。）

\section*{3.4\quad 拟凸函数}

\subsection*{3.4.1\quad 定义及例子}

函数 $f:\mathbb{R}^n\to\mathbb{R}$ 称为拟凸函数（或者单峰函数），如果其定义域及所有下水平集
\[
S_\alpha=\{x\in\dom f\mid f(x)\leq\alpha\},
\]
$\alpha\in\mathbb{R}$，都是凸集。函数 $f$ 是拟凹函数，如果 $-f$ 是拟凸函数，即每个上水平集 $\{x\mid f(x)\geq\alpha\}$ 是凸集。若某函数既是拟凸函数又是拟凹函数，其为拟线性函数。函数是拟线性函数，如果其定义域和所有的水平集 $\{x\mid f(x)=\alpha\}$ 都是凸集。

对于定义在 $\mathbb{R}$ 上的函数，拟凸性要求每个下水平集是一个区间（有可能包括无限区间）。$\mathbb{R}$ 上的一个拟凸函数如图 3.9 所示。

凸函数具有凸的下水平集，所以也是拟凸函数。但是拟凸函数不一定是凸函数，图 3.9 所示的简单例子即说明了这一点。

\noindent\textbf{例 3.29}\quad $\mathbb{R}$ 上的一些例子：
\begin{itemize}
\item 对数函数。定义在 $\mathbb{R}_{++}$ 上的函数 $\log x$ 是拟凸函数（也是拟凹函数，因此是拟线性函数）。
\item 上取整函数。函数 $\operatorname{ceil}(x)=\inf\{z\in\mathbb{Z}\mid z\geq x\}$ 是拟凸函数（亦为拟凹函数）。
\end{itemize}

从上述例子可以看出，拟凸函数可能是凹函数，甚至有可能是不连续的。下面给出 $\mathbb{R}^n$ 上的一些例子。

\noindent\textbf{例 3.30\quad 向量的长度。}\quad 定义 $x\in\mathbb{R}^n$ 的长度为非零分量的下标的最大值，即
\[
f(x)=\max\{i\mid x_i\neq0\}.
\]
（定义零向量的长度为零。）由于此函数的下水平集是子空间
\[
f(x)\leq\alpha\Longleftrightarrow x_i=0\quad \forall i=\lfloor\alpha\rfloor+1,\ldots,n.
\]
所以它在 $\mathbb{R}^n$ 上是拟凸函数。

\noindent\textbf{例 3.31}\quad 考虑函数 $f:\mathbb{R}^2\to\mathbb{R}$，其定义域为 $\dom f=\mathbb{R}_+^2$，$f(x_1,x_2)=x_1x_2$。此函数既非凸函数，亦非凹函数，因为其 Hessian 矩阵
\[
\nabla^2 f(x)=
\begin{bmatrix}
0 & 1\\
1 & 0
\end{bmatrix}
\]
是不定的：它的两个特征值一个大于零一个小于零。然而，函数 $f$ 是拟凹函数，因为对任意 $\alpha$，函数的上水平集
\[
\{x\in\mathbb{R}_+^2\mid x_1x_2\geq\alpha\}
\]
都是凸集。（注意到函数 $f$ 在 $\mathbb{R}^2$ 上不是拟凸函数。）

\noindent\textbf{例 3.32\quad 线性分式函数。}\quad 函数
\[
f(x)=\frac{a^Tx+b}{c^Tx+d},
\]
其定义域为 $\dom f=\{x\mid c^Tx+d>0\}$，是拟凸函数，也是拟凹函数，所以此函数是拟线性函数。其 $\alpha$-下水平集为
\[
\begin{aligned}
S_\alpha
&=\{x\mid c^Tx+d>0,\ (a^Tx+b)/(c^Tx+d)\leq\alpha\}\\
&=\{x\mid c^Tx+d>0,\ a^Tx+b\leq\alpha(c^Tx+d)\},
\end{aligned}
\]
它是凸集，因为它是一个开的半平面和闭的半平面的交集。（可以用同样的方法来判断其上水平集是凸集。）

\noindent\textbf{例 3.33\quad 距离比函数。}\quad 设 $a,b\in\mathbb{R}^n$，定义
\[
f(x)=\frac{\|x-a\|_2}{\|x-b\|_2},
\]
即，到 $a$ 点的 Euclid 距离和到 $b$ 点的 Euclid 距离之比。函数 $f$ 在半平面 $\{x\mid \|x-a\|_2\leq\|x-b\|_2\}$ 上是拟凸函数。为了说明这一点，我们考虑 $f$ 的 $\alpha$-下水平集，由于在半平面 $\{x\mid \|x-a\|_2\leq\|x-b\|_2\}$ 上 $f(x)\leq1$，选取 $\alpha\leq1$。下水平集是满足下式的一系列点
\[
\|x-a\|_2\leq\alpha\|x-b\|_2.
\]
将上式两端平方，并重新排列各项，我们得到如下表达式
\[
(1-\alpha^2)x^Tx-2(a-\alpha^2b)^Tx+a^Ta-\alpha^2b^Tb\leq0.
\]
当 $\alpha\leq1$ 时其为凸集（事实上是一个 Euclid 球）。

\noindent\textbf{例 3.34\quad 内生回报率。}\quad 令 $x=(x_0,x_1,\ldots,x_n)$ 表示 $n$ 个时间段内的现金流序列，$x_i>0$ 表明在时间段 $i$ 流入现金 $x_i$，$x_i<0$ 表明在时间段 $i$ 流出现金 $-x_i$。令利率 $r\geq0$，定义现金流的现值为
\[
\operatorname{PV}(x,r)=\sum_{i=0}^n(1+r)^{-i}x_i.
\]
（$(1+r)^{-i}$ 是时间段 $i$ 流入或者流出现金的折扣因子。）

我们考虑 $x_0<0$ 且 $x_0+x_1+\cdots+x_n>0$ 的现金流。这种情况表明初始投资 $|x_0|$，总的现金余额 $x_1+\cdots+x_n$ 多于最初的投资额（没有考虑任何折扣因子）。

对于上述现金流，$\operatorname{PV}(x,0)>0$，当 $r\to\infty$ 时 $\operatorname{PV}(x,r)\to x_0<0$，因此至少存在一个 $r\geq0$，使得 $\operatorname{PV}(x,r)=0$。定义现金流的内生回报率为使得现值为零的最小利率 $r\geq0$，即
\[
\operatorname{IRR}(x)=\inf\{r\geq0\mid \operatorname{PV}(x,r)=0\}.
\]

内生回报率是 $x$ 的拟凹函数（当 $x_0<0$ 且 $x_1+\cdots+x_n>0$ 时）。为了说明这一点，利用如下事实
\[
\operatorname{IRR}(x)\geq R\Longleftrightarrow \operatorname{PV}(x,r)\geq0\quad \forall 0\leq r<R.
\]
上式的左边定义了 $\operatorname{IRR}$ 的 $R$-上水平集。右端是关于不同的 $r$（$0\leq r<R$）的一族集合 $\{x\mid \operatorname{PV}(x,r)>0\}$ 的交集。对于任意 $r$，$\operatorname{PV}(x,r)>0$ 定义了一个开的半平面，所以上式的右端定义了一个凸集。

\subsection*{3.4.2\quad 基本性质}

上面的例子说明了拟凸性是凸性的重要扩展。在拟凸条件下，凸函数的很多性质仍然成立，或者可以找到类似性质。例如，存在一种变化的 Jensen 不等式来描述拟凸函数：函数 $f$ 是拟凸函数的充要条件是，$\dom f$ 是凸集，且对于任意 $x,y\in\dom f$ 及 $0\leq\theta\leq1$，有
\begin{equation}
f(\theta x+(1-\theta)y)\leq\max\{f(x),f(y)\};
\tag{3.19}
\end{equation}
即线段中任意一点的函数值不超过其端点函数值中最大的那个。不等式 (3.19) 有时称为拟凸函数的 Jensen 不等式，图 3.10 所示即为一个拟凸函数的例子。

\noindent\textbf{例 3.35\quad 非零向量的基数。}\quad 向量 $x\in\mathbb{R}^n$ 的基数或者规模定义为其中非零分量的个数，用 $\operatorname{card}(x)$ 表示。函数 $\operatorname{card}$ 在 $\mathbb{R}_+^n$ 上是拟凹函数（注意到不是 $\mathbb{R}^n$ 上）。由修正的 Jensen 不等式可以立即得到，对 $x,y\succeq0$，有
\[
\operatorname{card}(x+y)\geq\min\{\operatorname{card}(x),\operatorname{card}(y)\}.
\]

\noindent\textbf{例 3.36\quad 半正定矩阵的秩。}\quad 函数 $\operatorname{rank}X$ 在 $\mathbb{S}_+^n$ 上是拟凹函数。由修正的 Jensen 不等式 (3.19) 可以得到，对 $X,Y\in\mathbb{S}_+^n$ 下式成立
\[
\operatorname{rank}(X+Y)\geq\min\{\operatorname{rank}X,\operatorname{rank}Y\},
\]
（因为当 $x\succeq0$ 时，$\operatorname{rank}(\operatorname{diag}(x))=\operatorname{card}(x)$，此例可以看成是上一个例子的扩展。）

和凸性类似，拟凸性可以由函数 $f$ 在直线上的性质刻画：函数 $f$ 是拟凸的充要条件是它在和其定义域相交的任意直线上是拟凸函数。特别地，可以通过将一个函数限制在任意直线上，通过考察所得到的函数在 $\mathbb{R}$ 上的拟凸性来验证原函数的拟凸性。

\noindent\textbf{$\mathbb{R}$ 上的拟凸函数}\quad 对 $\mathbb{R}$ 上的拟凸函数，我们给出一个简单的刻画。由于考虑一般的函数较为繁琐，所以我们考虑连续函数。连续函数 $f:\mathbb{R}\to\mathbb{R}$ 是拟凸的，当且仅当下述条件至少有一个成立。
\begin{itemize}
\item 函数 $f$ 是非减的；
\item 函数 $f$ 是非增的；
\item 存在一点 $c\in\dom f$，使得对于 $t\leq c$（且 $t\in\dom f$），$f$ 非增，对于 $t\geq c$（且 $t\in\dom f$），$f$ 非减。
\end{itemize}
点 $c$ 可以在 $f$ 的全局最小点中任选一个。图 3.11 描述了这样的情形。

\subsection*{3.4.3\quad 可微拟凸函数}

\noindent\textbf{一阶条件}\quad 设函数 $f:\mathbb{R}^n\to\mathbb{R}$ 可微，则函数 $f$ 是拟凸的充要条件是，$\dom f$ 是凸集，且对于任意 $x,y\in\dom f$ 有
\begin{equation}
f(y)\leq f(x)\Longrightarrow \nabla f(x)^T(y-x)\leq0.
\tag{3.20}
\end{equation}
这是和不等式 (3.2) 相对应的描述拟凸函数的不等式。

当 $\nabla f(x)\neq0$ 时，条件 (3.20) 的几何意义很简单，即 $\nabla f(x)$ 在点 $x$ 处定义了水平集 $\{y\mid f(y)\leq f(x)\}$ 的一个支撑超平面，如图 3.12 所示。

我们注意到判断凸性的一阶条件 (3.2) 和判断拟凸性的一阶条件 (3.20) 相似，但是实际上二者存在重要的差别。例如，如果函数 $f$ 是凸函数且 $\nabla f(x)=0$，$x$ 是函数 $f$ 的全局极小点。然而，对于拟凸函数，这样的论断并不成立：有可能 $\nabla f(x)=0$，但是点 $x$ 不是 $f$ 的全局极小点。

\noindent\textbf{二阶条件}\quad 假设函数 $f$ 二次可微。如果函数 $f$ 是拟凸函数，则对于任意 $x\in\dom f$ 以及任意 $y\in\mathbb{R}^n$ 有
\begin{equation}
y^T\nabla f(x)=0\Longrightarrow y^T\nabla^2f(x)y\geq0.
\tag{3.21}
\end{equation}
对于定义在 $\mathbb{R}$ 上的拟凸函数，上述条件可以简化为如下条件
\[
f'(x)=0\Longrightarrow f''(x)\geq0,
\]
即在斜率为零的点，二阶导数是非负的。对于定义在 $\mathbb{R}^n$ 上的拟凸函数，条件 (3.21) 的意义就稍微复杂一些。和 $n=1$ 时的情形类似，只要 $\nabla f(x)=0$，有 $\nabla^2 f(x)\succeq0$。若 $\nabla f(x)\neq0$，条件 (3.21) 说明了 $\nabla^2f(x)$ 在 $(n-1)$-维子空间 $\nabla f(x)^\perp$ 上是半正定的。即 $\nabla^2f(x)$ 最多只有一个负特征值。

反之（部分条件），如果对于任意 $x\in\dom f$ 以及任意 $y\in\mathbb{R}^n$，函数 $f$ 满足
\begin{equation}
y^T\nabla f(x)=0\Longrightarrow y^T\nabla^2f(x)y>0
\tag{3.22}
\end{equation}
则函数 $f$ 是拟凸函数。此条件等价于在满足 $\nabla f(x)=0$ 的点处 $\nabla^2f(x)$ 正定，在其他点处 $\nabla^2f(x)$ 在 $(n-1)$-维子空间 $\nabla f(x)^\perp$ 上正定。

\noindent\textbf{拟凸性的二阶条件的证明}\quad 将函数限制在任意直线上，我们只要考虑 $f:\mathbb{R}\to\mathbb{R}$ 的情形即可判断函数的拟凸性。

首先我们说明，如果函数 $f:\mathbb{R}\to\mathbb{R}$ 在区间 $(a,b)$ 上是拟凸的，则其必须满足 (3.21)，即对 $c\in(a,b)$，若 $f'(c)=0$，则有 $f''(c)\geq0$。否则，若存在 $c\in(a,b)$，使得 $f'(c)=0$，但 $f''(c)<0$，则对于小的正数 $\epsilon$，有 $f(c-\epsilon)<f(c)$ 和 $f(c+\epsilon)<f(c)$。因此，对于小的正数 $\epsilon$，水平集 $\{x\mid f(x)\leq f(c)-\epsilon\}$ 是不连通的，也即不是凸集，这与 $f$ 是拟凸函数的假设相矛盾。

接下来我们说明如果条件 (3.22) 成立，则函数 $f$ 是拟凸的。假设式 (3.22) 成立，即对于区间 $c\in(a,b)$ 上满足 $f'(c)=0$ 的点，我们有 $f''(c)>0$。这意味着当函数 $f'$ 过零点时，它是严格增的。因此，函数 $f'$ 最多过一次零点。如果函数 $f'$ 不过零点，则 $f$ 在 $(a,b)$ 上或者非增，或者非减，因此是拟凸函数。若函数 $f'$ 恰过一次零点，设在点 $c\in(a,b)$ 处过零点。因为 $f''(c)>0$，则在 $a<t<c$ 时 $f'(t)\leq0$，在 $c\leq t<b$ 时，$f'(t)\geq0$。因此函数 $f$ 是拟凸的。

\subsection*{3.4.4\quad 保拟凸运算}

\noindent\textbf{非负加权最大}\quad 拟凸函数的非负加权最大定义为
\[
f=\max\{w_1f_1,\ldots,w_mf_m\},
\]
其中 $w_i\geq0$，$f_i$ 是拟凸函数。上述定义的函数 $f$ 是拟凸函数。此性质可以扩展到一般的逐点上确界，即
\[
f(x)=\sup_{y\in C}\left(w(y)g(x,y)\right),
\]
其中 $w(y)\geq0$。固定任意 $y$，$g(x,y)$ 关于 $x$ 是拟凸函数。上述定义的函数 $f$ 是拟凸函数。可以很容易证明上述结论：$f(x)\leq\alpha$，当且仅当
\[
w(y)g(x,y)\leq\alpha\quad \forall y\in C,
\]
即函数 $f$ 的 $\alpha$-下水平集是以 $x$ 为自变量的一族函数 $w(y)g(x,y)$ 的 $\alpha$-下水平集的交集。

\noindent\textbf{例 3.37\quad 广义特征值。}\quad 一对对称矩阵 $(X,Y)$，其中 $Y\succ0$，定义其最大广义特征值为
\[
\lambda_{\max}(X,Y)=\sup_{u\neq0}\frac{u^TXu}{u^TYu}
=\sup\{\lambda\mid \det(\lambda Y-X)=0\}.
\]
（参见 \S A.5.3。）此函数在定义域 $\dom f=\mathbb{S}^n\times\mathbb{S}_{++}^n$ 上是拟凸函数。

为了说明这一点，我们考虑如下表达式
\[
\lambda_{\max}(X,Y)=\sup_{u\neq0}\frac{u^TXu}{u^TYu}.
\]
对于任意 $u\neq0$，函数 $u^TXu/u^TYu$ 是 $(X,Y)$ 的线性分式函数，因此是 $(X,Y)$ 的拟凸函数。因为 $\lambda_{\max}$ 是一族拟凸函数的上确界，因此它是拟凸函数。

\noindent\textbf{复合}\quad 如果函数 $g:\mathbb{R}^n\to\mathbb{R}$ 是拟凸函数，且函数 $h:\mathbb{R}\to\mathbb{R}$ 是非减的，则复合函数 $f=h\circ g$ 是拟凸函数。

拟凸函数和一个仿射函数或者线性分式函数进行复合可以得到拟凸函数。如果函数 $f$ 是拟凸函数，则 $g(x)=f(Ax+b)$ 是拟凸函数，且函数 $\tilde g(x)=f((Ax+b)/(c^Tx+d))$ 在集合
\[
\{x\mid c^Tx+d>0,\ (Ax+b)/(c^Tx+d)\in\dom f\}
\]
上是拟凸函数。

\noindent\textbf{最小化}\quad 如果函数 $f(x,y)$ 是 $x$ 和 $y$ 的联合拟凸函数，且 $C$ 是凸集，则函数
\[
g(x)=\inf_{y\in C}f(x,y)
\]
是拟凸函数。

为了说明这一点，我们只需要证明对任意 $\alpha\in\mathbb{R}$，集合 $\{x\mid g(x)\leq\alpha\}$ 是凸集。根据函数 $g$ 的定义，$g(x)\leq\alpha$，当且仅当对任意 $\epsilon>0$，存在一 $y\in C$ 使得 $f(x,y)\leq\alpha+\epsilon$。

令 $x_1$ 和 $x_2$ 是 $g$ 的 $\alpha$-下水平集中的任意两点。则对于任意 $\epsilon>0$，存在 $y_1,y_2\in C$ 使得
\[
f(x_1,y_1)\leq\alpha+\epsilon,\quad
f(x_2,y_2)\leq\alpha+\epsilon,
\]
因为函数 $f$ 关于 $x$ 和 $y$ 是拟凸的，我们有
\[
f(\theta x_1+(1-\theta)x_2,\theta y_1+(1-\theta)y_2)\leq\alpha+\epsilon,
\]
其中 $0\leq\theta\leq1$。因此 $g(\theta x_1+(1-\theta)x_2)\leq\alpha$，说明 $\{x\mid g(x)\leq\alpha\}$ 是凸集。

\subsection*{3.4.5\quad 通过一族凸函数进行表示}

下面我们将要说明，可以很方便地将拟凸函数 $f$ 的下水平集（凸集）表示成凸函数的不等式。选择一族凸函数 $\phi_t:\mathbb{R}^n\to\mathbb{R}$，$t\in\mathbb{R}$ 表示凸函数的编号，这些函数满足
\begin{equation}
f(x)\leq t\Longleftrightarrow \phi_t(x)\leq0,
\tag{3.23}
\end{equation}
即拟凸函数 $f$ 的 $t$-下水平集是凸函数 $\phi_t$ 的 $0$-下水平集。显然，对于任意 $x\in\mathbb{R}^n$，函数 $\phi_t$ 必须满足：当 $s\geq t$ 时，$\phi_t(x)\leq0\Longrightarrow\phi_s(x)\leq0$。为了满足这个条件，要求对于任意 $x$，$\phi_t(x)$ 是 $t$ 的非增函数，即当 $s\geq t$ 时，$\phi_s(x)\leq\phi_t(x)$。

为了说明总能找到这样一族函数，我们可以选取
\[
\phi_t(x)=
\begin{cases}
0, & f(x)\leq t,\\
\infty, & \text{其他情况},
\end{cases}
\]
即 $\phi_t$ 是函数 $f$ 的 $t$-下水平集的示性函数。显然这样的一族函数不是唯一的，例如如果函数 $f$ 的下水平集是闭集，我们可以选取
\[
\phi_t(x)=\operatorname{dist}(x,\{z\mid f(z)\leq t\}).
\]
当然，我们希望选择的 $\phi_t$ 具有良好的性质，比如说可微性。

\noindent\textbf{例 3.38\quad 凸凹函数之比。}\quad 设 $p$ 是凸函数，$q$ 是凹函数，在凸集 $C$ 上，$p(x)\geq0$，$q(x)>0$。在集合 $C$ 上定义函数 $f(x)=p(x)/q(x)$，则 $f$ 是拟凸函数。

此时有
\[
f(x)\leq t\Longleftrightarrow p(x)-tq(x)\leq0,
\]
因此我们可以选取 $\phi_t(x)=p(x)-tq(x)$，$t\geq0$。对于任意 $t$，$\phi_t$ 是凸函数，且对于任意 $x$，$\phi_t(x)$ 是关于 $t$ 的减函数。

\section*{3.5\quad 对数-凹函数和对数-凸函数}

\subsection*{3.5.1\quad 定义}

称函数 $f:\mathbb{R}^n\to\mathbb{R}$ 对数凹或对数-凹，如果对所有的 $x\in\dom f$ 有 $f(x)>0$ 且 $\log f$ 是凹函数。称函数对数凸或者对数-凸，如果 $\log f$ 是凸函数。因此，函数 $f$ 是对数-凸的，当且仅当 $1/f$ 是对数-凹的。也可以允许函数值 $f$ 为零，这时只需在 $f$ 为零的点令 $\log f(x)=-\infty$。此时，函数 $f$ 是对数-凹的，如果扩展值函数 $\log f$ 是凹函数。

对数-凹的性质可以不借助对数直接表达：考察函数 $f:\mathbb{R}^n\to\mathbb{R}$，其定义域是凸集，且对于任意 $x\in\dom f$ 有 $f(x)>0$，函数是对数-凹的，当且仅当对任意 $x,y\in\dom f$，$0\leq\theta\leq1$，有
\[
f(\theta x+(1-\theta)y)\geq f(x)^\theta f(y)^{1-\theta}.
\]
特别地，对数-凹函数在两点之间中点的函数值不小于这两点的函数值的几何平均值。

根据函数复合规则，我们知道如果函数 $h$ 是凸函数，则函数 $e^h$ 是凸函数，因此对数-凸函数是凸函数。类似地，非负凹函数是对数-凹函数。此外，由于对数函数是单调增函数，所以对数-凸函数是拟凸函数，对数-凹函数是拟凹函数。

\noindent\textbf{例 3.39\quad 一些简单的对数-凹函数和对数-凸函数。}
\begin{itemize}
\item 仿射函数。函数 $f(x)=a^Tx+b$ 在 $\{x\mid a^Tx+b>0\}$ 上是对数-凹函数。
\item 幂函数。函数 $f(x)=x^a$ 在 $\mathbb{R}_{++}$ 上当 $a\leq0$ 时是对数-凸函数，当 $a\geq0$ 时是对数-凹函数。
\item 指数函数。函数 $f(x)=e^{ax}$ 既是对数-凸函数也是对数-凹函数。
\item Gauss 概率密度函数的累积分布函数
\[
\Phi(x)=\frac{1}{\sqrt{2\pi}}\int_{-\infty}^x e^{-u^2/2}\,du,
\]
是对数-凹函数（参见习题 3.54）。
\item Gamma 函数。Gamma 函数
\[
\Gamma(x)=\int_0^\infty u^{x-1}e^{-u}\,du,
\]
当 $x\geq1$ 时是对数-凸函数（参见习题 3.52）。
\item 行列式。$\det X$ 在 $\mathbb{S}_{++}^n$ 上是对数-凹函数。
\item 行列式与迹之比。$\det X/\operatorname{tr}X$ 在 $\mathbb{S}_{++}^n$ 上是对数-凹函数（参见习题 3.49）。
\end{itemize}

\noindent\textbf{例 3.40\quad 对数-凹的概率密度函数。}\quad 一些常用的概率密度函数是对数-凹函数。其中的两个例子是多变量正态分布的概率密度函数，即
\[
f(x)=\frac{1}{\sqrt{(2\pi)^n\det\Sigma}}
e^{-\frac{1}{2}(x-\bar x)^T\Sigma^{-1}(x-\bar x)}
\]
（其中 $\bar x\in\mathbb{R}^n$，$\Sigma\in\mathbb{S}_{++}^n$），以及 $\mathbb{R}_+^n$ 上的指数分布的概率密度函数
\[
f(x)=\left(\prod_{i=1}^n\lambda_i\right)e^{-\lambda^Tx}
\]
（其中 $\lambda\succ0$）。另外还有一个例子是凸集 $C$ 上均匀分布的概率密度函数
\[
f(x)=
\begin{cases}
1/\alpha, & x\in C,\\
0, & x\notin C,
\end{cases}
\]
其中 $\alpha=\operatorname{vol}(C)$ 表示集合 $C$ 的体积（Lebesgue 测度）。在这样的定义下，若 $x$ 在集合 $C$ 外，$\log f$ 取值为 $-\infty$，若 $x$ 在集合 $C$ 内，$\log f$ 取值为 $-\log\alpha$，是凹函数。

我们考虑一个不太常见的分布，Wishart 分布，其定义如下。令 $x_1,\ldots,x_p\in\mathbb{R}^n$ 是独立的 Gauss 随机向量，均值为零，协方差为 $\Sigma\in\mathbb{S}_{++}^n$，其中 $p>n$。随机矩阵 $X=\sum_{i=1}^p x_ix_i^T$ 具有 Wishart 概率密度函数
\[
f(X)=a(\det X)^{(p-n-1)/2}e^{-\frac{1}{2}\operatorname{tr}(\Sigma^{-1}X)},
\]
其中 $\dom f=\mathbb{S}_{++}^n$，$a$ 是正常数。Wishart 概率密度函数是对数-凹函数，这是因为
\[
\log f(X)=\log a+\frac{p-n-1}{2}\log\det X-\frac{1}{2}\operatorname{tr}(\Sigma^{-1}X)
\]
是 $X$ 的凹函数。

\subsection*{3.5.2\quad 相关性质}

\noindent\textbf{二次可微的对数-凸/凹函数}\quad 设函数 $f$ 二次可微，其中 $\dom f$ 是凸集，我们有
\[
\nabla^2\log f(x)=\frac{1}{f(x)}\nabla^2 f(x)-\frac{1}{f(x)^2}\nabla f(x)\nabla f(x)^T.
\]
事实上，函数 $f$ 是对数-凸函数，当且仅当对任意 $x\in\dom f$，下式成立
\[
f(x)\nabla^2 f(x)\succeq\nabla f(x)\nabla f(x)^T,
\]
函数 $f$ 是对数-凹函数，当且仅当对任意 $x\in\dom f$，下式成立
\[
f(x)\nabla^2 f(x)\preceq\nabla f(x)\nabla f(x)^T.
\]

\noindent\textbf{乘积，和，以及积分运算}\quad 对数-凸性以及对数-凹性对乘积以及正的伸缩运算是封闭的。例如，如果函数 $f$ 和 $g$ 是对数-凹函数，则逐点乘积函数 $h(x)=f(x)g(x)$ 也是对数-凹函数，这是因为 $\log h(x)=\log f(x)+\log g(x)$，而 $\log f(x)$ 和 $\log g(x)$ 是 $x$ 的凹函数。

一些简单的例子即可以说明对数-凹函数的和一般不是对数-凹函数。然而，对数-凸函数的和仍然是对数-凸函数。设 $f$ 和 $g$ 是对数-凸函数，即 $F=\log f$ 和 $G=\log g$ 是凸函数。根据凸函数的复合规则有
\[
\log(\exp F+\exp G)=\log(f+g);
\]
上述函数是凸函数。因此，两个对数-凸函数的和仍然是对数-凸函数。

更一般地，如果对任意 $y\in C$，$f(x,y)$ 是 $x$ 的对数-凸函数，则函数
\[
g(x)=\int_C f(x,y)\,dy
\]
是对数-凸函数。

\noindent\textbf{例 3.41\quad 非负函数的 Laplace 变换以及矩生成函数和累积量生成函数。}\quad 设函数 $p:\mathbb{R}^n\to\mathbb{R}$ 对所有 $x$ 满足 $p(x)\geq0$。函数 $p$ 的 Laplace 变换
\[
P(z)=\int p(x)e^{-z^Tx}\,dx,
\]
在 $\mathbb{R}^n$ 上是对数-凸函数。（此时定义域 $\dom P$ 为 $\{z\mid P(z)<\infty\}$。）

假设 $p$ 是概率密度函数，即满足 $\int p(x)\,dx=1$。函数 $M(z)=P(-z)$ 称为概率密度函数的矩生成函数。之所以称为矩生成函数是因为对 $M(z)$ 求导并令 $z=0$ 可以得到概率密度函数的矩，即
\[
\nabla M(0)=\mathbf{E}v,\quad
\nabla^2M(0)=\mathbf{E}vv^T,
\]
其中随机向量 $v$ 具有概率密度函数 $p$。

凸函数 $\log M(z)$ 被称为函数 $p$ 的累积量生成函数，这是因为通过它的导数可以得到概率密度函数的累积量。例如，对其求一阶以及二阶导，并令自变量为零，可以得到相应随机变量的期望和协方差为
\[
\nabla\log M(0)=\mathbf{E}v,\quad
\nabla^2\log M(0)=\mathbf{E}(v-\mathbf{E}v)(v-\mathbf{E}v)^T.
\]

\noindent\textbf{对数-凹函数的积分}\quad 在一些特殊的情况下，对数-凹的性质在积分后仍然保留。如果函数 $f:\mathbb{R}^n\times\mathbb{R}^m\to\mathbb{R}$ 是对数-凹函数，则
\[
g(x)=\int f(x,y)\,dy
\]
在 $\mathbb{R}^n$ 上是 $x$ 的对数-凹函数。（此时是在 $\mathbb{R}^m$ 上求积分。）这个结论的证明不太简单，参见参考文献。

这个结论有着重要的意义，本节后续将会就其中的一部分进行描述。根据此结论可知对数-凹的概率密度分布函数的边际分布仍然是对数-凹的。此外，对数-凹的性质对卷积运算是封闭的，即如果函数 $f$ 和 $g$ 在 $\mathbb{R}^n$ 上是对数-凹的，则它们的卷积
\[
(f*g)(x)=\int f(x-y)g(y)\,dy
\]
仍然是对数-凹函数。（这个不难看出，因为 $g(y)$ 和 $f(x-y)$ 关于 $(x,y)$ 是对数-凹的，因此乘积 $f(x-y)g(y)$ 是对数-凹函数，根据前面关于积分的结论就可以直接得到此处的结论。）

设 $C\subseteq\mathbb{R}^n$ 是凸集，$w$ 是 $\mathbb{R}^n$ 上的随机向量，设其具有对数-凹的概率密度函数 $p$。则函数
\[
f(x)=\operatorname{prob}(x+w\in C)
\]
是 $x$ 的对数-凹函数。为了说明这一点，我们将 $f$ 表述为
\[
f(x)=\int g(x+w)p(w)\,dw,
\]
其中 $g$ 定义为
\[
g(u)=
\begin{cases}
1, & u\in C,\\
0, & u\notin C,
\end{cases}
\]
（它是对数-凹函数），利用前面关于积分的结论即可得到 $f$ 是对数-凹函数。

\noindent\textbf{例 3.42\quad 概率密度函数的累积分布函数。}\quad $f:\mathbb{R}^n\to\mathbb{R}$ 定义为
\[
F(x)=\operatorname{prob}(w\leq x)=\int_{-\infty}^{x_n}\cdots\int_{-\infty}^{x_1} f(z)\,dz_1\cdots dz_n,
\]
其中 $w$ 是随机变量，具有概率密度函数 $f$。如果函数 $f$ 是对数-凹函数，那么 $F$ 是对数-凹的。之前我们已经接触过一个这样的例子，Gauss 随机变量的累积分布函数
\[
f(x)=\frac{1}{\sqrt{2\pi}}\int_{-\infty}^{x}e^{-t^2/2}\,dt,
\]
是对数-凹函数。（参见例 3.39 和习题 3.54。）

\noindent\textbf{例 3.43\quad 产出函数。}\quad 令 $x\in\mathbb{R}^n$ 表示生产产品的一组参数的名义值或者目标值。不同制造过程会使得参数值有一定变化，得到的值为 $x+w$，其中 $w\in\mathbb{R}^n$ 是随机向量，表征了制造过程的变化，其均值为零。制造过程的产出是名义参数值的函数，
\[
Y(x)=\operatorname{prob}(x+w\in S),
\]
其中 $S\subseteq\mathbb{R}^n$ 表示可以接受的产品参数值的集合，即产品规格。

如果制造误差 $w$ 的概率密度函数是对数-凹的（例如 Gauss 分布），产品规格集合 $S$ 是凸集，那么产出函数 $Y$ 是对数-凹函数。定义 $\alpha$-产出区域为使产出超过 $\alpha$ 的名义参数值的集合，由于 $Y$ 是对数-凹函数，此集合是凸集。例如，$95\%$ 产出区域
\[
\{x\mid Y(x)\geq0.95\}=\{x\mid \log Y(x)\geq \log0.95\}
\]
是凸集，这是因为它是凹函数 $\log Y$ 的上水平集。

\noindent\textbf{例 3.44\quad 多面体的体积。}\quad 令 $A\in\mathbb{R}^{m\times n}$，定义
\[
P_u=\{x\in\mathbb{R}^n\mid Ax\preceq u\}.
\]
则其体积 $\operatorname{vol}P_u$ 是 $u$ 的对数-凹函数。

为了证明这一点，注意到函数
\[
\Psi(x,u)=
\begin{cases}
1, & Ax\preceq u,\\
0, & \text{其他情况},
\end{cases}
\]
是对数-凹的。根据前面关于积分的结论，我们有
\[
\int\Psi(x,u)\,dx=\operatorname{vol}P_u
\]
是对数-凹的。

\section*{3.6\quad 关于广义不等式的凸性}

本节考虑广义的单调性和凸性，采用广义不等式，而不是通常的 $\mathbb{R}$ 上的顺序。

\subsection*{3.6.1\quad 关于广义不等式的单调性}

设 $K\subseteq\mathbb{R}^n$ 是一个正常锥，其相应的广义不等式为 $\preceq_K$。称函数 $f:\mathbb{R}^n\to\mathbb{R}$ $K$-非减，如果下式成立
\[
x\preceq_K y\Longrightarrow f(x)\leq f(y),
\]
称函数 $K$-增，如果下式成立
\[
x\preceq_K y,\ x\neq y\Longrightarrow f(x)<f(y).
\]
类似地，可以定义 $K$-非增和 $K$-减函数。

\noindent\textbf{例 3.45\quad 单调向量函数。}\quad 函数 $f:\mathbb{R}^n\to\mathbb{R}$ 在 $\mathbb{R}_+^n$ 上非减，当且仅当对任意 $x,y$ 下式成立
\[
x_1\leq y_1,\ldots,x_n\leq y_n\Longrightarrow f(x)\leq f(y).
\]
这等价于函数 $f$ 限制为每个分量 $x_i$ 的函数时（$x_i$ 作为自变量，固定 $x_j$，$j\neq i$）是非减的。

\noindent\textbf{例 3.46\quad 矩阵单调函数。}\quad 称函数 $f:\mathbb{S}^n\to\mathbb{R}$ 矩阵单调（增或减），如果在半正定锥内函数是单调的。下面给出一些定义在 $X\in\mathbb{S}^n$ 上的矩阵单调函数的例子：
\begin{itemize}
\item 函数 $\operatorname{tr}(WX)$，其中 $W\in\mathbb{S}^n$，当 $W\succeq0$ 时，函数是矩阵非减的，当 $W\succ0$ 时是矩阵增的（当 $W\preceq0$ 时是矩阵非增，当 $W\prec0$ 时是矩阵减）。
\item 函数 $\operatorname{tr}(X^{-1})$ 在 $\mathbb{S}_{++}^n$ 上矩阵减。
\item 函数 $\det X$ 在 $\mathbb{S}_{++}^n$ 上矩阵增，在 $\mathbb{S}_+^n$ 上矩阵非减。
\end{itemize}

\noindent\textbf{单调性的梯度条件}\quad 在前面的章节中，我们提到，可微函数 $f:\mathbb{R}\to\mathbb{R}$，其定义域是凸集（即一个区间），函数是非减的，当且仅当对任意 $x\in\dom f$ 有 $f'(x)\geq0$，函数是增的，如果对任意 $x\in\dom f$ 有 $f'(x)>0$（反过来不一定成立）。在本节中，可以很方便地将这些条件扩展至关于广义不等式的单调性。考虑可微函数 $f$，其定义域为凸集，它是 $K$-非减的，当且仅当对任意 $x\in\dom f$ 有
\begin{equation}
\nabla f(x)\succeq_{K^*}0.
\tag{3.24}
\end{equation}
注意到和简单的标量情况的区别：梯度在对偶不等式中必须是非负的。对于严格的情况，我们有，如果对任意 $x\in\dom f$ 存在
\begin{equation}
\nabla f(x)\succ_{K^*}0,
\tag{3.25}
\end{equation}
则函数 $f$ 是 $K$-增的。和标量情况类似，对于严格的情形，反过来不一定正确。

下面证明单调性的一阶条件。首先，假设对任意 $x$，函数 $f$ 满足式 (3.24)，但是它不是 $K$-非减的，即存在 $x,y$，$x\preceq_K y$ 且 $f(y)<f(x)$。由 $f$ 是可微的知道存在某个 $t\in[0,1]$ 使得下式成立
\[
\frac{d}{dt}f(x+t(y-x))=\nabla f(x+t(y-x))^T(y-x)<0.
\]
因为 $y-x\in K$，上式意味着
\[
\nabla f(x+t(y-x))\notin K^*,
\]
这和式 (3.24) 在每一点都成立的假设矛盾。类似地，我们可以证明若式 (3.25) 成立，则函数 $f$ 是 $K$-增的。

可以很容易地看出如果函数 $f$ 是 $K$-非减的，定义域中的每一点都必须满足式 (3.24)。否则，假设在 $x=z$ 处式 (3.24) 不成立。由对偶锥的定义，存在某一 $v\in K$ 使得
\[
\nabla f(z)^Tv<0.
\]
考虑 $t$ 的函数 $h(t)=f(z+tv)$。其导数 $h'(0)=\nabla f(z)^Tv<0$，因此存在 $t>0$ 使得 $h(t)=f(z+tv)<h(0)=f(z)$，即函数 $f$ 不是 $K$-非减的，导出矛盾。

\subsection*{3.6.2\quad 关于广义不等式的凸性}

设 $K\subseteq\mathbb{R}^m$ 是一正常锥，相应的广义不等式为 $\preceq_K$。函数 $f:\mathbb{R}^n\to\mathbb{R}^m$ 是 $K$-凸的，如果对于任意 $x,y$，以及 $0\leq\theta\leq1$ 有
\[
f(\theta x+(1-\theta)y)\preceq_K \theta f(x)+(1-\theta)f(y).
\]
函数被称为严格 $K$-凸的，如果对任意 $x\neq y$ 和 $0<\theta<1$ 有
\[
f(\theta x+(1-\theta)y)\prec_K \theta f(x)+(1-\theta)f(y).
\]
当 $m=1$ 时（即 $K=\mathbb{R}_+$），上述定义即为常见的凸以及严格凸的定义。

\noindent\textbf{例 3.47\quad 关于分量不等式的凸性。}\quad 函数 $f:\mathbb{R}^n\to\mathbb{R}^m$ 关于分量不等式（即 $K=\mathbb{R}_+^m$ 时的广义不等式）是凸的，当且仅当对任意 $x,y$ 和 $0\leq\theta\leq1$，有
\[
f(\theta x+(1-\theta)y)\preceq \theta f(x)+(1-\theta)f(y),
\]
即每个分量 $f_i$ 都是凸函数。函数 $f$ 关于分量不等式是严格凸的，当且仅当每个分量 $f_i$ 是严格凸的。

\noindent\textbf{例 3.48\quad 矩阵凸性。}\quad 设函数 $f$ 是对称矩阵值函数，即 $f:\mathbb{R}^n\to\mathbb{S}^m$。称函数 $f$ 关于矩阵不等式是凸的，如果对任意 $x$ 和 $y$ 以及 $\theta\in[0,1]$，有
\[
f(\theta x+(1-\theta)y)\preceq \theta f(x)+(1-\theta)f(y).
\]
这种凸性有时称为矩阵凸性。一个等价的定义是对任意向量 $z$，标量函数 $z^Tf(x)z$ 都是凸函数。（这是一个证明矩阵凸性的好方法。）称矩阵函数为严格矩阵凸的，如果对 $x\neq y$ 和 $0<\theta<1$，有
\[
f(\theta x+(1-\theta)y)\prec \theta f(x)+(1-\theta)f(y),
\]
或者等价地，如果对任意 $z\neq0$ 函数 $z^Tfz$ 严格凸。

一些例子。
\begin{itemize}
\item 函数 $f(X)=XX^T$，其中 $X\in\mathbb{R}^{n\times m}$，是矩阵凸的，这是因为任选 $z$，函数 $z^TXX^Tz=\|X^Tz\|^2$ 是 $X$（分量）的凸的二次函数。同理，函数 $f(X)=X^2$ 在 $\mathbb{S}^n$ 上也是矩阵凸的。
\item 当 $1\leq p\leq2$ 或 $-1\leq p\leq0$ 时，函数 $X^p$ 在 $\mathbb{S}_{++}^n$ 上是矩阵凸的，当 $0\leq p\leq1$ 时，函数是矩阵凹的。
\item 当 $n\geq2$ 时，函数 $f(X)=e^X$ 在 $\mathbb{S}^n$ 上不是矩阵凸的。
\end{itemize}

凸函数的很多结论都可以扩展到 $K$-凸函数。我们给出一个简单的例子，例如，函数是 $K$-凸的，当且仅当它在定义域上的任意直线上是 $K$-凸的。本节后续部分会提到一些关于 $K$-凸的结论，我们将会在后面用到这些结论；更多的结论可以参看习题。

\noindent\textbf{$K$-凸的对偶刻画}\quad 函数 $f$ 是 $K$-凸的，当且仅当对任意 $w\succeq_{K^*}0$，（实值）函数 $w^Tf$ 是凸的（常规凸性）；$f$ 是严格 $K$-凸的，当且仅当对任意非零向量 $w\succeq_{K^*}0$，函数 $w^Tf$ 是严格凸的。（从对偶不等式的定义和性质可以直接得到上述结论。）

\noindent\textbf{可微的 $K$-凸函数}\quad 可微函数 $f$ 是 $K$-凸的，当且仅当其定义域是凸集，且对任意 $x,y\in\dom f$ 有
\[
f(y)\succeq_K f(x)+Df(x)(y-x).
\]
（其中 $Df(x)\in\mathbb{R}^{m\times n}$ 是函数 $f$ 在 $x$ 处的导数或者 Jacobian 矩阵；参看 \S A.4.1。）称函数 $f$ 是严格 $K$-凸的，当且仅当对任意 $x,y\in\dom f$，$x\neq y$，有
\[
f(y)\succ_K f(x)+Df(x)(y-x).
\]

\noindent\textbf{复合定理}\quad 函数复合保留凸性的很多结论都可以推广到 $K$-凸的情形。例如，如果函数 $g:\mathbb{R}^n\to\mathbb{R}^p$ 是 $K$-凸的，函数 $h:\mathbb{R}^p\to\mathbb{R}$ 是凸的，$\tilde h$（$h$ 的扩展值延伸）是 $K$-非减的，那么函数 $h\circ g$ 是凸的。回忆之前的一般凸性的情形，以凸函数为自变量的非减凸函数仍然是凸函数，此处的结论是对这个结论的推广。前提条件 $\tilde h$ 是 $K$-非减的意味着
\[
\dom h-K=\dom h.
\]

\noindent\textbf{例 3.49\quad 定义二次矩阵函数 $g:\mathbb{R}^{m\times n}\to\mathbb{S}^n$}
\[
g(X)=X^TAX+B^TX+X^TB+C,
\]
其中 $A\in\mathbb{S}^m$，$B\in\mathbb{R}^{m\times n}$ 以及 $C\in\mathbb{S}^n$。当 $A\succeq0$ 时，函数是凸的。

函数 $h:\mathbb{S}^n\to\mathbb{R}$，$h(Y)=-\log\det(-Y)$，在 $\dom h=-\mathbb{S}_{++}^n$ 上是凸且增的。

根据复合定理，我们有，函数
\[
f(X)=-\log\det(-(X^TAX+B^TX+X^TB+C))
\]
在
\[
\dom f=\{X\in\mathbb{R}^{m\times n}\mid X^TAX+B^TX+X^TB+C\prec0\}
\]
上是凸的。这是对下面的结论的推广，即当 $a>0$ 时，函数
\[
-\log(-(ax^2+bx+c))
\]
在
\[
\{x\in\mathbb{R}\mid ax^2+bx+c<0\}
\]
上是凸的。

\section*{参考文献}

凸分析的标准参考文献是 Rockafellar [Roc70]。其他关于凸函数的书有 Stoer 和 Witzgall [SW70]，Roberts 和 Varberg [RV73]，Van Tiel [vT84]，Hiriart-Urruty 和 Lemaréchal [HUL93]，Ekeland 和 Témam [ET99]，Borwein 和 Lewis [BL00]，Florenzano 和 Le Van [FL01]，Barvinok [Bar02]，以及 Bertsekas，Nedić 和 Ozdaglar [Ber03]。很多非线性规划的教材也有一些章节涉及凸函数（例如，Mangasarian [Man94]，Bazaraa，Sherali 和 Shetty [BSS93]，Bertsekas [Ber99]，Polyak [Pol87]，以及 Peressini，Sullivan 和 Uhl [PSU88]）。

文献 [Jen06] 中提到了 Jensen 不等式。在不等式的一般性研究中，Jensen 不等式起到了重要的作用，Hardy，Littlewood 和 Pólya [HLP52] 以及 Beckenbach 和 Bellman [BB65] 中都涉及了不等式的研究。

透视函数的概念来源于 Hiriart-Urruty 和 Lemaréchal [HUL93，第 1 册，第 100 页]。例 3.19 中的定义（相对熵和 Kullback-Leibler 散度），以及相应的习题 3.13，可以参看 Cover 和 Thomas [CT91]。

早期的一些关于拟凸函数（以及凸性的其他扩展）的重要参考文献有 Nikaido [Nik54]，Mangasarian [Man94，第 9 章]，Arrow 和 Enthoven [AE61]，Ponstein [Pon67]，以及 Luenberger [Lue68]。更全面的此类参考文献可以参照 Bazaraa，Sherali 和 Shetty [BSS93，第 126 页]。Prékopa [Pré80] 对对数-凹函数做了一个综述。在 Barndorff-Nielsen [BN78，\S7] 中提到了 Laplace 变换的对数-凸性。关于对数-凹函数的积分的结论证明可以参看 Prékopa [Pré71, Pré73]。

广义不等式在最近的关于锥优化的参考文献中广泛使用，如 Nesterov 和 Nemirovski [NN94，第 156 页]；Ben-Tal 和 Nemirovski [BTN01]以及第 4 章最后所列的参考文献。关于广义不等式的凸性在 Luenberger [Lue69，\S8.2] 和 Isii [Isi64] 中亦有涉及。矩阵单调性和矩阵凸性由 Löwner [Löw34] 提出，Davis [Dav63]，Roberts 和 Varberg [RV73，第 216 页] 以及 Marshall 和 Olkin [MO79，\S16E] 对其进行了详细讨论。例 3.48 中提到的函数 $X^p$ 的凸性或者凹性的相关结论可以参看 Bondar [Bon94，定理 16.1]。另一个简单的例子，函数 $e^X$ 不是矩阵凸的，可以参看 Marshall 和 Olkin [MO79，第 474 页]。

\section*{习题}

\noindent\textbf{凸性的定义}

\noindent\textbf{3.1}\quad 假设 $f:\mathbb{R}\to\mathbb{R}$ 是凸函数，$a,b\in\dom f$，$a<b$。

(a) 证明对任意 $x\in[a,b]$，下式成立
\[
f(x)\leq \frac{b-x}{b-a}f(a)+\frac{x-a}{b-a}f(b).
\]

(b) 证明对任意 $x\in(a,b)$，下式成立
\[
\frac{f(x)-f(a)}{x-a}\leq \frac{f(b)-f(a)}{b-a}\leq \frac{f(b)-f(x)}{b-x}.
\]
画一个草图来描述此不等式。

(c) 假设函数 $f$ 可微。利用 (b) 中的结果来证明
\[
f'(a)\leq \frac{f(b)-f(a)}{b-a}\leq f'(b).
\]
注意到这些不等式也可以通过式 (3.2) 得到：
\[
f(b)\geq f(a)+f'(a)(b-a),\qquad
f(a)\geq f(b)+f'(b)(a-b).
\]

(d) 假设函数 $f$ 二次可微。利用 (c) 中的结论证明 $f''(a)\geq0$ 以及 $f''(b)\geq0$。

\noindent\textbf{3.2\quad 凸，凹，拟凸以及拟凹函数的水平集。}\quad 函数 $f$ 的一些水平集如下面左图所示。曲线 1 表示 $\{x\mid f(x)=1\}$，其他曲线类似。

函数 $f$ 是否可以为凸函数（凹函数，拟凸函数，拟凹函数）？解释你的答案。对下面右图的曲线进行同样的判断。

\noindent\textbf{3.3\quad 一个增的凸函数的反函数。}\quad 设 $f:\mathbb{R}\to\mathbb{R}$ 递增，在其定义域 $(a,b)$ 上是凸函数。令 $g$ 表示其反函数，即具有定义域 $(f(a),f(b))$，且对所有 $a<x<b$ 满足 $g(f(x))=x$。函数 $g$ 是凸函数还是凹函数？为什么？

\noindent\textbf{3.4}\quad [RV73，第 15 页] 证明连续函数 $f:\mathbb{R}^n\to\mathbb{R}$ 是凸函数的充要条件是，对任意线段，函数在线段上的平均值不大于线段端点函数值的平均：即对任意 $x,y\in\mathbb{R}^n$，下式成立
\[
\int_0^1 f(x+\lambda(y-x))\,d\lambda\leq \frac{f(x)+f(y)}{2}.
\]

\noindent\textbf{3.5}\quad [RV73，第 22 页] \textbf{凸函数的滑动平均。} 设函数 $f:\mathbb{R}\to\mathbb{R}$ 是凸函数，$\mathbb{R}_+\subseteq\dom f$。证明其滑动平均 $F$，即
\[
F(x)=\frac{1}{x}\int_0^x f(t)\,dt,\qquad \dom F=\mathbb{R}_{++},
\]
是凸函数。提示：固定 $s$，$f(sx)$ 是 $x$ 的凸函数，所以 $\int_0^1 f(sx)\,ds$ 是凸函数。

\noindent\textbf{3.6\quad 函数及上境图。}\quad 什么时候函数的上境图是半平面？什么时候函数的上境图是一个凸锥？什么时候函数的上境图是一个多面体？

\noindent\textbf{3.7}\quad 设函数 $f:\mathbb{R}^n\to\mathbb{R}$ 是凸函数，其定义域为 $\dom f=\mathbb{R}^n$，函数在 $\mathbb{R}^n$ 上有上界。证明函数 $f$ 是常数。

\noindent\textbf{3.8\quad 凸性的二阶条件。}\quad 证明二次可微函数 $f$ 是凸函数的充要条件是，其定义域是凸集，且对任意 $x\in\dom f$ 有 $\nabla^2 f(x)\succeq0$。提示：首先考虑 $f:\mathbb{R}\to\mathbb{R}$ 的情形。可以利用凸性的一阶条件（在第 63 页已经证明）。

\noindent\textbf{3.9\quad 仿射集上的凸性的二阶条件。}\quad 设 $F\in\mathbb{R}^{n\times m}$，$\hat x\in\mathbb{R}^n$。函数 $f:\mathbb{R}^n\to\mathbb{R}$ 限制在仿射集 $\{Fz+\hat x\mid z\in\mathbb{R}^m\}$ 上的函数定义为 $\tilde f:\mathbb{R}^m\to\mathbb{R}$，其满足
\[
\tilde f(z)=f(Fz+\hat x),\qquad
\dom\tilde f=\{z\mid Fz+\hat x\in\dom f\}.
\]
设函数 $f$ 的定义域是凸集，且在定义域上函数二次可微。

(a) 证明函数 $\tilde f$ 是凸函数的充要条件是，对任意 $z\in\dom\tilde f$，有
\[
F^T\nabla^2 f(Fz+\hat x)F\succeq0.
\]

(b) 设 $A\in\mathbb{R}^{p\times n}$ 是一矩阵，其零空间为矩阵 $F$ 的值域（矩阵列向量张成的线性空间），即 $AF=0$ 且 $\operatorname{rank}A=n-\operatorname{rank}F$。证明函数 $\tilde f$ 是凸函数的充要条件是，对任意 $z\in\dom\tilde f$，存在一 $\lambda\in\mathbb{R}$ 使得下式成立
\[
\nabla^2 f(Fz+\hat x)+\lambda A^TA\succeq0.
\]

提示：利用如下结论：如果 $B\in\mathbb{S}^n$ 以及 $A\in\mathbb{R}^{p\times n}$，则对任意 $x\in\mathcal{N}(A)$ 式 $x^TBx\geq0$ 都成立的充要条件是，存在一 $\lambda$ 使得 $B+\lambda A^TA\succeq0$。

\noindent\textbf{3.10\quad Jensen 不等式的一个扩展。}\quad 关于 Jensen 不等式的一个理解是随机化或者波动不利，即会提高一个凸函数的平均值：若函数 $f$ 是凸函数，$v$ 是随机变量，均值为零，我们有 $\mathbf{E}f(x_0+v)\geq f(x_0)$。我们于是有下面的猜想。如果 $f_0$ 是凸函数，那么 $v$ 的方差越大，期望 $\mathbf{E}f(x_0+v)$ 越大。

(a) 给出一个反例，说明这样的猜想是不正确的。即找到两个均值为零的随机变量 $v$ 和 $w$，$\operatorname{var}(v)>\operatorname{var}(w)$，一个凸函数 $f$ 以及一点 $x_0$，使得 $\mathbf{E}f(x_0+v)<\mathbf{E}f(x_0+w)$。

(b) 当 $v$ and $w$ 之间具有比例关系时上面的猜想是正确的。给定凸函数 $f$ 以及零均值随机变量 $v$，证明当 $t\geq0$ 时 $\mathbf{E}f(x_0+tv)$ 关于 $t$ 单调增加。

\noindent\textbf{3.11\quad 单调映射。}\quad 称函数 $\psi:\mathbb{R}^n\to\mathbb{R}^n$ 是单调的，如果对任意 $x,y\in\dom\psi$，下式成立
\[
(\psi(x)-\psi(y))^T(x-y)\geq0.
\]
（注意到此处的“单调性”不同于 \S3.6.1 给出的定义。但是两种定义都被广泛使用。）设函数 $f:\mathbb{R}^n\to\mathbb{R}$ 是一可微凸函数。证明其梯度 $\nabla f$ 是单调的。反过来，是否每个单调映射都对应某凸函数的梯度？

\noindent\textbf{3.12}\quad 设函数 $f:\mathbb{R}^n\to\mathbb{R}$ 是凸函数，$g:\mathbb{R}^n\to\mathbb{R}$ 是凹函数，$\dom f=\dom g=\mathbb{R}^n$，对任意 $x$ 有 $g(x)\leq f(x)$。证明存在一个仿射函数 $h$，使得对于任意 $x$，$g(x)\leq h(x)\leq f(x)$。换言之，如果凹函数 $g$ 是凸函数 $f$ 的一个下估计，那么我们可以在函数 $f$ 和 $g$ 之间找到一个仿射函数。

\noindent\textbf{3.13\quad Kullback-Leibler 散度和信息不等式。}\quad 令 $D_{\mathrm{kl}}$ 表示 Kullback-Leibler 散度，如式 (3.17) 所定义。证明信息不等式：对任意 $u,v\in\mathbb{R}_{++}^n$，有 $D_{\mathrm{kl}}(u,v)\geq0$。此外，证明当且仅当 $u=v$ 时，$D_{\mathrm{kl}}(u,v)=0$。

提示：Kullback-Leibler 散度可以表示为
\[
D_{\mathrm{kl}}(u,v)=f(u)-f(v)-\nabla f(v)^T(u-v);
\]
其中 $f(v)=\sum_{i=1}^n v_i\log v_i$ 是 $v$ 的负熵。

\noindent\textbf{3.14\quad 凸-凹函数和鞍点。}\quad 我们说函数 $f:\mathbb{R}^n\times\mathbb{R}^m\to\mathbb{R}$ 是凸-凹函数，如果任意固定 $x$，$f(x,z)$ 是 $z$ 的凹函数，任意固定 $z$，$f(x,z)$ 是 $x$ 的凸函数。同时，要求函数的定义域具有积的形式，$\dom f=A\times B$，其中 $A\subseteq\mathbb{R}^n$ 和 $B\subseteq\mathbb{R}^m$ 都是凸集。

(a) 对二次可微函数 $f:\mathbb{R}^n\times\mathbb{R}^m\to\mathbb{R}$，从 Hessian 矩阵 $\nabla^2 f(x,z)$ 的角度给出函数为凸-凹函数的二阶条件。

(b) 设函数 $f:\mathbb{R}^n\times\mathbb{R}^m\to\mathbb{R}$ 是凸-凹函数并且可微，$\nabla f(\bar x,\bar z)=0$。证明鞍点性质成立：即对任意 $x,z$，有
\[
f(\bar x,z)\leq f(\bar x,\bar z)\leq f(x,\bar z).
\]
证明上述性质成立可以推导出函数 $f$ 满足强极大极小性质：
\[
\sup_z\inf_x f(x,z)=\inf_x\sup_z f(x,z)
\]
（且等式两端共同值为 $f(\bar x,\bar z)$）。

(c) 设函数 $f:\mathbb{R}^n\times\mathbb{R}^m\to\mathbb{R}$ 是可微的，但不一定是凸-凹函数，若在点 $\bar x,\bar z$ 处鞍点性质成立：即对任意 $x,z$，有
\[
f(\bar x,z)\leq f(\bar x,\bar z)\leq f(x,\bar z).
\]
证明 $\nabla f(\bar x,\bar z)=0$。

\noindent\textbf{例子}

\noindent\textbf{3.15\quad 一族凹的效用函数。}\quad 对 $0<\alpha\leq1$，令
\[
u_\alpha(x)=\frac{x^\alpha-1}{\alpha},
\]
其中 $\dom u_\alpha=\mathbb{R}_+$。我们定义 $u_0(x)=\log x$（其定义域为 $\dom u_0=\mathbb{R}_{++}$）。

(a) 证明对 $x>0$，$u_0(x)=\lim_{\alpha\to0}u_\alpha(x)$。

(b) 证明 $u_\alpha$ 是凹函数，单调增加，且都满足 $u_\alpha(1)=0$。

上述函数在经济学中经常被用来对一定数量的商品的收益或者资金的效用进行建模。$u_\alpha$ 是凹函数意味着边际效用（即商品增加固定的数量时所带来的收益的增加）随着商品数量的增加而减少。换言之，凹性反映了饱和效果。

\noindent\textbf{3.16}\quad 判断下列函数是否是凸函数，凹函数，拟凸函数以及拟凹函数？

(a) 函数 $f(x)=e^x-1$，定义域为 $\mathbb{R}$。

(b) 函数 $f(x_1,x_2)=x_1x_2$，定义域为 $\mathbb{R}_{++}^2$。

(c) 函数 $f(x_1,x_2)=1/(x_1x_2)$，定义域为 $\mathbb{R}_{++}^2$。

(d) 函数 $f(x_1,x_2)=x_1/x_2$，定义域为 $\mathbb{R}_{++}^2$。

(e) 函数 $f(x_1,x_2)=x_1^2/x_2$，定义域为 $\mathbb{R}\times\mathbb{R}_{++}$。

(f) 函数 $f(x_1,x_2)=x_1^\alpha x_2^{1-\alpha}$，其中 $0\leq\alpha\leq1$，定义域为 $\mathbb{R}_{++}^2$。

\noindent\textbf{3.17}\quad 设 $p<1$，$p\neq0$。证明定义域为 $\dom f=\mathbb{R}_{++}^n$ 的函数
\[
f(x)=\left(\sum_{i=1}^n x_i^p\right)^{1/p}
\]
是凹函数。这其中包括两类特殊的函数 $f(x)=\left(\sum_{i=1}^n x_i^{1/2}\right)^2$ 以及调和平均函数 $f(x)=\left(\sum_{i=1}^n 1/x_i\right)^{-1}$。提示：仿照 \S3.1.5 中关于指数函数和的对数以及几何平均函数的证明并进行适当修改。

\noindent\textbf{3.18}\quad 仿照 \S3.1.5 中对数-行列式函数凹性的证明并进行适当修改来证明下面的命题。

(a) 函数 $f(X)=\operatorname{tr}(X^{-1})$ 在 $\dom f=\mathbb{S}_{++}^n$ 上是凸函数。

(b) 函数 $f(X)=(\det X)^{1/n}$ 在 $\dom f=\mathbb{S}_{++}^n$ 上是凹函数。

\noindent\textbf{3.19\quad 非负加权和以及积分。}

(a) 证明函数 $f(x)=\sum_{i=1}^r\alpha_i x_{[i]}$ 是关于 $x$ 的凸函数，其中 $\alpha_1\geq\alpha_2\geq\cdots\geq\alpha_r\geq0$，$x_{[i]}$ 表示 $x$ 第 $i$ 大的分量。（可以利用 $f(x)=\sum_{i=1}^k x_{[i]}$ 在 $\mathbb{R}^n$ 上是凸函数来进行证明。）

(b) 令 $T(x,\omega)$ 表示三角多项式
\[
T(x,\omega)=x_1+x_2\cos\omega+x_3\cos2\omega+\cdots+x_n\cos(n-1)\omega.
\]
证明函数
\[
f(x)=-\int_0^{2\pi}\log T(x,\omega)\,d\omega
\]
在 $\{x\in\mathbb{R}^n\mid T(x,\omega)>0,\ 0\leq\omega\leq2\pi\}$ 上是凸函数。

\noindent\textbf{3.20\quad 与仿射函数的复合。}\quad 证明下列函数 $f:\mathbb{R}^n\to\mathbb{R}$ 是凸函数。

(a) 函数 $f(x)=\|Ax-b\|$，其中 $A\in\mathbb{R}^{m\times n}$，$b\in\mathbb{R}^m$，$\|\cdot\|$ 是 $\mathbb{R}^m$ 上的范数。

(b) 函数 $f(x)=-(\det(A_0+x_1A_1+\cdots+x_nA_n))^{1/m}$，其定义域为 $\{x\mid A_0+x_1A_1+\cdots+x_nA_n\succ0\}$，其中 $A_i\in\mathbb{S}^m$。

(c) 函数 $f(X)=\operatorname{tr}(A_0+x_1A_1+\cdots+x_nA_n)^{-1}$，其定义域为 $\{x\mid A_0+x_1A_1+\cdots+x_nA_n\succ0\}$，其中 $A_i\in\mathbb{S}^m$。（利用 $\operatorname{tr}(X^{-1})$ 在 $\mathbb{S}_{++}^n$ 上是凸函数的结论来进行证明；参见习题 3.18。）

\noindent\textbf{3.21\quad 逐点最大和上确界。}\quad 证明下列函数 $f:\mathbb{R}^n\to\mathbb{R}$ 是凸函数。

(a) 函数 $f(x)=\max_{i=1,\ldots,k}\|A^{(i)}x-b^{(i)}\|$，其中 $A^{(i)}\in\mathbb{R}^{m\times n}$，$b^{(i)}\in\mathbb{R}^m$，$\|\cdot\|$ 是 $\mathbb{R}^m$ 上的范数。

(b) 函数 $f(x)=\sum_{i=1}^r |x|_{[i]}$，其定义域为 $\mathbb{R}^n$，其中 $|x|$ 表示向量，其分量 $|x|_i=|x_i|$（即 $|x|$ 的每个分量是向量 $x$ 对应分量的绝对值），$|x|_{[i]}$ 表示向量 $|x|$ 中第 $i$ 大的分量。换言之，$|x|_{[1]},|x|_{[2]},\ldots,|x|_{[n]}$ 是向量 $x$ 的分量的绝对值按照非增的顺序进行排列。

\noindent\textbf{3.22\quad 复合规则。}\quad 证明下列函数是凸函数。

(a) 函数
\[
f(x)=-\log\left(-\log\left(\sum_{i=1}^m e^{a_i^Tx+b_i}\right)\right),
\]
其定义域为
\[
\dom f=\left\{x\mid \sum_{i=1}^m e^{a_i^Tx+b_i}<1\right\}.
\]
可以利用函数 $\log\left(\sum_{i=1}^n e^{y_i}\right)$ 是凸函数的结论来进行证明。

(b) 函数 $f(x,u,v)=-\sqrt{uv-x^Tx}$，其定义域为
\[
\dom f=\{(x,u,v)\mid uv>x^Tx,\ u,v>0\}.
\]
可以应用如下事实来进行证明：函数 $x^Tx/u$ 当 $u>0$ 时是 $(x,u)$ 的凸函数；函数 $-\sqrt{x_1x_2}$ 在 $\mathbb{R}_{++}^2$ 上是凸函数。

(c) 函数 $f(x,u,v)=-\log(uv-x^Tx)$，其定义域为
\[
\dom f=\{(x,u,v)\mid uv>x^Tx,\ u,v>0\}.
\]

(d) 函数 $f(x,t)=-(t^p-\|x\|_p^p)^{1/p}$，其中 $p>1$，定义域为
\[
\dom f=\{(x,t)\mid t\geq\|x\|_p\}.
\]
可以应用下面两个结论来进行证明：即 $\|x\|_p^p/u^{p-1}$ 在 $u>0$ 时是 $(x,u)$ 的凸函数（参见习题 3.23）以及 $-x^{1/p}y^{1-1/p}$ 在 $\mathbb{R}_{++}^2$ 上是凸函数（参见习题 3.16）。

(e) 函数 $f(x,t)=-\log(t^p-\|x\|_p^p)$，其中 $p>1$，定义域为
\[
\dom f=\{(x,t)\mid t>\|x\|_p\}.
\]
可以应用 $\|x\|_p^p/u^{p-1}$ 当 $u>0$ 时是 $(x,u)$ 的凸函数这个结论来进行证明（参见习题 3.23）。

\noindent\textbf{3.23\quad 函数的透视函数。}

(a) 证明当 $p>1$ 时，函数
\[
f(x,t)=\frac{|x_1|^p+\cdots+|x_n|^p}{t^{p-1}}=\frac{\|x\|_p^p}{t^{p-1}}
\]
在 $\{(x,t)\mid t>0\}$ 上是凸函数。

(b) 证明函数
\[
f(x)=\frac{\|Ax+b\|_2^2}{c^Tx+d}
\]
在 $\{x\mid c^Tx+d>0\}$ 上是凸函数，其中 $A\in\mathbb{R}^{m\times n}$，$b\in\mathbb{R}^m$，$c\in\mathbb{R}^n$ 以及 $d\in\mathbb{R}$。

\noindent\textbf{3.24\quad 概率单纯形上的一些函数。}\quad 令 $x$ 是实值随机变量，其取值的集合为 $\{a_1,\ldots,a_n\}$，其中 $a_1<a_2<\cdots<a_n$，$\operatorname{prob}(x=a_i)=p_i$，$i=1,\ldots,n$。对下列 $p$ 的函数（在概率单纯形 $\{p\in\mathbb{R}_+^n\mid \mathbf{1}^Tp=1\}$ 上），判断函数是否为凸函数，凹函数，拟凸函数或者拟凹函数。

(a) 函数 $\mathbf{E}x$。

(b) 函数 $\operatorname{prob}(x\geq\alpha)$。

(c) 函数 $\operatorname{prob}(\alpha\leq x\leq\beta)$。

(d) 函数 $\sum_{i=1}^n p_i\log p_i$，即分布的负熵函数。

(e) 函数 $\operatorname{var}x=\mathbf{E}(x-\mathbf{E}x)^2$。

(f) 函数 $\operatorname{quartile}(x)=\inf\{\beta\mid \operatorname{prob}(x\leq\beta)\geq0.25\}$。

(g) 以 $>90\%$ 的概率落入的最小集合 $A\subseteq\{a_1,\ldots,a_n\}$ 的基数。（此处的基数意味着集合 $A$ 中元素的个数。）

(h) 以 $90\%$ 概率落入的最小区间，即
\[
\inf\{|\beta-\alpha|\mid \operatorname{prob}(\alpha\leq x\leq\beta)\geq0.9\}.
\]

\noindent\textbf{3.25\quad 不同分布间的最大概率距离。}\quad 令 $p,q\in\mathbb{R}^n$ 表示 $\{1,\ldots,n\}$ 上的两个概率分布，（因此 $p,q\succeq0$，$\mathbf{1}^Tp=\mathbf{1}^Tq=1$）。对所有事件，在概率分布 $p$ 和 $q$ 下发生概率的差异的最大值定义为概率分布 $p$ 和 $q$ 之间的最大概率距离 $d_{\mathrm{mp}}(p,q)$，即
\[
d_{\mathrm{mp}}(p,q)=\max\{|\operatorname{prob}(p,C)-\operatorname{prob}(q,C)|\mid C\subseteq\{1,\ldots,n\}\}.
\]
其中 $\operatorname{prob}(p,C)$ 是在概率分布 $p$ 下事件 $C$ 发生的概率，即 $\operatorname{prob}(p,C)=\sum_{i\in C}p_i$。

利用 $\|p-q\|_1=\sum_{i=1}^n |p_i-q_i|$ 将 $d_{\mathrm{mp}}$ 的表达式进行简化，并证明函数 $d_{\mathrm{mp}}$ 是 $\mathbb{R}^n\times\mathbb{R}^n$ 上的凸函数。（其定义域为 $\{(p,q)\mid p,q\succeq0,\ \mathbf{1}^Tp=\mathbf{1}^Tq=1\}$，但是可以很容易将其扩展到全空间 $\mathbb{R}^n\times\mathbb{R}^n$。）

\noindent\textbf{3.26\quad 更多关于特征值的函数。}\quad 令 $\lambda_1(X)\geq\lambda_2(X)\geq\cdots\geq\lambda_n(X)$ 表示矩阵 $X\in\mathbb{S}^n$ 的特征值。之前已经讨论过一些有关特征值的函数，它们是 $X$ 的凸函数或者凹函数。
\begin{itemize}
\item 最大特征值函数 $\lambda_1(X)$ 是凸函数（见例 3.10）。最小特征值函数 $\lambda_n(X)$ 是凹函数。
\item 特征值的和（或者矩阵的迹）$\operatorname{tr}X=\lambda_1(X)+\cdots+\lambda_n(X)$，是线性函数。
\item 特征值倒数的和（或者逆矩阵的迹），$\operatorname{tr}(X^{-1})=\sum_{i=1}^n 1/\lambda_i(X)$，是 $\mathbb{S}_{++}^n$ 上的凸函数（见习题 3.18）。
\item 特征值的几何平均 $(\det X)^{1/n}=\left(\prod_{i=1}^n\lambda_i(X)\right)^{1/n}$，以及特征值乘积的对数 $\log\det X=\sum_{i=1}^n\log\lambda_i(X)$，在 $X\in\mathbb{S}_{++}^n$ 上是凹函数（见习题 3.18 以及第 68 页的描述）。
\end{itemize}

在这道题中，利用变分特征，我们讨论更多的关于特征值的函数。

(a) 最大 $k$ 个特征值的和。证明函数 $\sum_{i=1}^k\lambda_i(X)$ 在 $\mathbb{S}^n$ 上是凸函数。提示：[HJ85，第 191 页]利用变分特征
\[
\sum_{i=1}^k\lambda_i(X)=\sup\{\operatorname{tr}(V^TXV)\mid V\in\mathbb{R}^{n\times k},\ V^TV=I\}.
\]

(b) 最小 $k$ 个特征值的几何平均。证明函数
\[
\left(\prod_{i=n-k+1}^n\lambda_i(X)\right)^{1/k}
\]
在 $\mathbb{S}_{++}^n$ 上是凹函数。提示：[MO79，第 513 页]对 $X\succ0$，我们有
\[
\left(\prod_{i=n-k+1}^n\lambda_i(X)\right)^{1/k}
=\frac{1}{k}\inf\{\operatorname{tr}(V^TXV)\mid V\in\mathbb{R}^{n\times k},\ \det V^TV=1\}.
\]

(c) 最小 $k$ 个特征值乘积的对数。证明函数 $\sum_{i=n-k+1}^n\log\lambda_i(X)$ 在 $\mathbb{S}_{++}^n$ 上是凹函数。提示：[MO79，第 513 页]对 $X\succ0$，有
\[
\prod_{i=n-k+1}^n\lambda_i(X)
=\inf\left\{\prod_{i=1}^k (V^TXV)_{ii}\ \middle|\ V\in\mathbb{R}^{n\times k},\ V^TV=I\right\}.
\]

\noindent\textbf{3.27\quad Cholesky 分解的因子的对角线元素。}\quad 任意矩阵 $X\in\mathbb{S}_{++}^n$ 有着唯一的 Cholesky 分解 $X=LL^T$，其中 $L$ 是下三角矩阵，$L_{ii}>0$。证明函数 $L_{ii}$ 是 $X$ 的凹函数（其定义域为 $\mathbb{S}_{++}^n$）。

提示：函数 $L_{ii}$ 可以表示为 $L_{ii}=(w-z^TY^{-1}z)^{1/2}$，其中
\[
\begin{bmatrix}
Y & z\\
z^T & w
\end{bmatrix}
\]
是矩阵 $X$ 左上角 $i\times i$ 子矩阵。

\noindent\textbf{保凸运算}

\noindent\textbf{3.28\quad 将凸函数表示为一系列仿射函数的逐点上确界。}\quad 在这道题中，我们将第 77 页证明的结论扩展至 $\dom f\neq\mathbb{R}^n$ 的情形。令 $f:\mathbb{R}^n\to\mathbb{R}$ 是凸函数，定义 $\bar f:\mathbb{R}^n\to\mathbb{R}$ 为 $f$ 的所有仿射全局下估计函数的逐点上确界：
\[
\bar f(x)=\sup\{g(x)\mid g\text{ 仿射},\ g(z)\leq f(z)\ \forall z\}.
\]

(a) 证明对 $x\in\operatorname{int}\dom f$ 有 $f(x)=\bar f(x)$。

(b) 证明如果 $f$ 是闭的（即 $\operatorname{epi} f$ 是闭集，见 \S A.3.3），有 $f=\bar f$。

\noindent\textbf{3.29\quad 分片线性凸函数的表示。}\quad 函数 $f:\mathbb{R}^n\to\mathbb{R}$，其定义域为 $\dom f=\mathbb{R}^n$，被称为分片线性函数，如果存在 $\mathbb{R}^n$ 的一个划分
\[
\mathbb{R}^n=X_1\cup X_2\cup\cdots\cup X_L,
\]
其中 $\operatorname{int}X_i\neq\emptyset$，且对任意 $i\neq j$，$\operatorname{int}X_i\cap\operatorname{int}X_j=\emptyset$；以及一系列仿射函数 $a_1^Tx+b_1,\ldots,a_L^Tx+b_L$ 使得对 $x\in X_i$ 有 $f(x)=a_i^Tx+b_i$。

证明若 $f$ 是凸函数有 $f(x)=\max\{a_1^Tx+b_1,\ldots,a_L^Tx+b_L\}$。

\noindent\textbf{3.30\quad 函数的凸包或凸包络。}\quad 函数 $f:\mathbb{R}^n\to\mathbb{R}$ 的凸包或凸包络定义为
\[
g(x)=\inf\{t\mid (x,t)\in\operatorname{conv}\operatorname{epi} f\}.
\]
几何上，函数 $g$ 的上境图是 $f$ 的上境图的凸包。

证明函数 $g$ 是 $f$ 最大的凸下估计函数。换言之，证明如果函数 $h$ 是凸函数，且对所有 $x$ 有 $h(x)\leq f(x)$，则对所有 $x$，有 $h(x)\leq g(x)$。

\noindent\textbf{3.31}\quad [Roc70，第 35 页] \textbf{最大的齐次下估计函数。} 设 $f$ 是凸函数。定义函数 $g$ 为
\[
g(x)=\inf_{\alpha>0}\frac{f(\alpha x)}{\alpha}.
\]

(a) 证明函数 $g$ 是齐次的（即对任意 $t\geq0$，$g(tx)=tg(x)$）。

(b) 证明 $g$ 是 $f$ 的最大齐次下估计函数：即如果函数 $h$ 是齐次的，且对所有 $x$ 有 $h(x)\leq f(x)$，则对任意 $x$，有 $h(x)\leq g(x)$。

(c) 证明 $g$ 是凸函数。

\noindent\textbf{3.32\quad 凸函数的积或比。}\quad 一般而言，两个凸函数的乘积或者比值不是凸函数。但是对定义在 $\mathbb{R}$ 上的凸函数有一些较好的结论。证明下列结论。

(a) 如果定义在某个区间上的函数 $f$ 和 $g$ 都是凸函数，且都非减（或者都非增），二者都大于零，则函数 $fg$ 在此区间上是凸函数。

(b) 如果函数 $f$ 和 $g$ 是凹函数，大于零，一个非减，另一个非增，那么函数 $fg$ 是凹函数。

(c) 如果函数 $f$ 是凸函数，非减且大于零，函数 $g$ 是凹函数，非增，且大于零，那么函数 $f/g$ 是凸函数。

\noindent\textbf{3.33\quad 透视定理的直接证明。}\quad 直接证明 \S3.2.6 中所定义的凸函数 $f$ 的透视函数 $g$ 是凸函数：即证明 $\dom g$ 是凸集，且对于 $(x,t),(y,s)\in\dom g$ 以及 $0\leq\theta\leq1$，有
\[
g(\theta x+(1-\theta)y,\theta t+(1-\theta)s)\leq \theta g(x,t)+(1-\theta)g(y,s).
\]

\noindent\textbf{3.34\quad Minkowski 函数。}\quad 凸集 $C$ 的 Minkowski 函数定义为
\[
M_C(x)=\inf\{t>0\mid t^{-1}x\in C\}.
\]

(a) 如何寻找 $M_C(x)$？画图给出几何解释。

(b) 证明函数 $M_C$ 是齐次函数，即对 $\alpha\geq0$ 有 $M_C(\alpha x)=\alpha M_C(x)$。

(c) 函数 $M_C$ 的定义域 $\dom M_C$ 是什么？

(d) 证明函数 $M_C$ 是凸函数。

(e) 设集合 $C$ 是闭集，对称（即如果 $x\in C$，有 $-x\in C$）且内部不为空。证明 $M_C$ 是一种范数。相应的单位球是什么？

\noindent\textbf{3.35\quad 支撑函数演算。}\quad 我们知道，集合 $C\subseteq\mathbb{R}^n$ 的支撑函数定义为 $S_C(y)=\sup\{y^Tx\mid x\in C\}$。在第 75 页我们曾证明 $S_C$ 是凸函数。

(a) 证明 $S_B=S_{\operatorname{conv}B}$。

(b) 证明 $S_{A+B}=S_A+S_B$。

(c) 证明 $S_{A\cup B}=\max\{S_A,S_B\}$。

(d) 设 $B$ 是闭凸集。证明 $A\subseteq B$ 的充要条件是，对任意 $y$ 有 $S_A(y)\leq S_B(y)$。

\noindent\textbf{共轭函数}

\noindent\textbf{3.36}\quad 推导下列函数的共轭函数。

(a) 最大值函数。函数 $f(x)=\max_{i=1,\ldots,n}x_i$，定义在 $\mathbb{R}^n$ 上。

(b) 最大若干分量的和。函数 $f(x)=\sum_{i=1}^r x_{[i]}$，定义域为 $\mathbb{R}^n$。

(c) 定义在 $\mathbb{R}$ 上的分片线性函数。定义在 $\mathbb{R}$ 上的分片线性函数 $f(x)=\max_{i=1,\ldots,m}(a_ix+b_i)$。在求解过程中，可以假设 $a_i$ 按升序排列，即 $a_1\leq\cdots\leq a_m$，且每个函数 $a_ix+b_i$ 都不是冗余的，即任选 $k$，至少存在一点 $x$ 使得 $f(x)=a_kx+b_k$。

(d) 幂函数。定义在 $\mathbb{R}_{++}$ 上的函数 $f(x)=x^p$，其中 $p>1$。如果 $p<0$ 呢？

(e) 几何平均。定义在 $\mathbb{R}_{++}^n$ 上的几何平均函数 $f(x)=-(\prod x_i)^{1/n}$。

(f) 二阶锥上的负广义对数。函数 $f(x,t)=-\log(t^2-x^Tx)$，定义域为 $\{(x,t)\in\mathbb{R}^n\times\mathbb{R}\mid \|x\|_2<t\}$。

\noindent\textbf{3.37}\quad 给定函数 $f(X)=\operatorname{tr}(X^{-1})$，其定义域为 $\dom f=\mathbb{S}_{++}^n$。证明 $f(X)$ 的共轭函数为
\[
f^*(Y)=-2\operatorname{tr}(-Y)^{1/2},\qquad \dom f^*=-\mathbb{S}_+^n.
\]
提示：函数 $f$ 的梯度为 $\nabla f(X)=-X^{-2}$。

\noindent\textbf{3.38\quad Young 不等式。}\quad 设 $f:\mathbb{R}\to\mathbb{R}$ 是增函数，$f(0)=0$，设函数 $g$ 是其反函数。定义 $F$ 和 $G$
\[
F(x)=\int_0^x f(a)\,da,\qquad
G(y)=\int_0^y g(a)\,da.
\]
证明函数 $F$ 和 $G$ 互为共轭函数。给出简单图示解释 Young 不等式
\[
xy\leq F(x)+G(y).
\]

\noindent\textbf{3.39\quad 共轭函数的性质。}

(a) 凸函数与仿射函数的和的共轭函数。设 $g(x)=f(x)+c^Tx+d$，其中，函数 $f$ 是凸函数。利用 $f$ 的共轭函数 $f^*$（和 $c,d$）表达 $g$ 的共轭函数 $g^*$。

(b) 透视函数的共轭函数。将凸函数 $f$ 的透视函数的共轭函数用 $f^*$ 来进行表达。

(c) 共轭以及极小化。设函数 $f(x,z)$ 是 $(x,z)$ 的凸函数，定义 $g(x)=\inf_z f(x,z)$。利用 $f^*$ 表达 $g^*$。

作为一个应用，定义函数 $g(x)=\inf_z\{h(z)\mid Az+b=x\}$，其中 $h$ 是凸函数，用 $h^*,A$ 以及 $b$ 表达 $g(x)$ 的共轭函数。

(d) 共轭的共轭。证明闭凸函数的共轭的共轭是它自己：即如果 $f$ 是闭凸函数，$f=f^{**}$。（函数称为闭的，如果其上境图是闭集；见 \S A.3.3。）提示：证明 $f^{**}$ 是函数 $f$ 的所有仿射全局下估计函数的逐点上确界。然后利用习题 3.28 中的结论。

\noindent\textbf{3.40\quad 共轭函数的梯度及 Hessian 矩阵。}\quad 设函数 $f:\mathbb{R}^n\to\mathbb{R}$ 是凸函数且连续二次可微。设 $\bar y$ 和 $\bar x$ 满足 $\bar y=\nabla f(\bar x)$，$\nabla^2 f(\bar x)\succ0$。

(a) 证明 $\nabla f^*(\bar y)=\bar x$。

(b) 证明 $\nabla^2 f^*(\bar y)=\nabla^2 f(\bar x)^{-1}$。

\noindent\textbf{3.41\quad 标准化的负熵函数的共轭。}\quad 给定标准化的负熵函数
\[
f(x)=\sum_{i=1}^n x_i\log\frac{x_i}{\mathbf{1}^Tx},
\]
其中 $\dom f=\mathbb{R}_+^n$。证明其共轭函数为
\[
f^*(y)=
\begin{cases}
0, & \sum_{i=1}^n e^{y_i}\leq1,\\
+\infty, & \text{其他情况}.
\end{cases}
\]

\noindent\textbf{拟凸函数}

\noindent\textbf{3.42\quad 逼近宽度。}\quad 设 $f_0,\ldots,f_n:\mathbb{R}\to\mathbb{R}$ 为连续函数，考虑用 $f_1,\ldots,f_n$ 的线性组合去逼近函数 $f_0$。对 $x\in\mathbb{R}^n$，我们称
\[
f=x_1f_1+\cdots+x_nf_n
\]
在区间 $[0,T]$ 内以容许误差 $\epsilon>0$ 逼近函数 $f_0$，如果对任意 $0\leq t\leq T$ 存在 $|f(t)-f_0(t)|\leq\epsilon$。固定误差容限 $\epsilon>0$，考虑 $f$ 在区间 $[0,T]$ 逼近 $f_0$，逼近宽度定义为在此误差容限下最长的区间宽度 $T$
\[
W(x)=\sup\{T\mid |x_1f_1(t)+\cdots+x_nf_n(t)-f_0(t)|\leq\epsilon,\quad \forall 0\leq t\leq T\}.
\]
证明 $W$ 是拟凹函数。

\noindent\textbf{3.43\quad 拟凸性的一阶条件。}\quad 证明 \S3.4.3 中给出的判断拟凸性的一阶条件：可微函数 $f:\mathbb{R}^n\to\mathbb{R}$，其定义域 $\dom f$ 是凸集，则函数 $f$ 是拟凸函数的充要条件是，对任意 $x,y\in\dom f$ 有
\[
f(y)\leq f(x)\Longrightarrow \nabla f(x)^T(y-x)\leq0.
\]
提示：可以先考虑定义在 $\mathbb{R}$ 上的函数；一般维空间的结论可以通过将函数限制在定义域内任意直线上得到。

\noindent\textbf{3.44\quad 拟凸性的二阶条件。}\quad \S3.4.3 给出了描述拟凸性的二阶条件。在本题中，我们给出二阶条件的另一种形式。证明如下结论。

(a) 点 $x\in\dom f$ 满足式 (3.21) 的充要条件是，存在 $\sigma$ 使得下式成立
\begin{equation}
\nabla^2 f(x)+\sigma\nabla f(x)\nabla f(x)^T\succeq0.
\tag{3.26}
\end{equation}
对任意 $y\neq0$，点 $x\in\dom f$ 满足式 (3.22) 的充要条件是，存在 $\sigma$ 使得下式成立
\begin{equation}
\nabla^2 f(x)+\sigma\nabla f(x)\nabla f(x)^T\succ0.
\tag{3.27}
\end{equation}
提示：不失一般性，我们可以假设 $\nabla^2 f(x)$ 是对角阵。

(b) 点 $x\in\dom f$ 满足式 (3.21) 的充要条件是，$\nabla f(x)=0$ 且 $\nabla^2 f(x)\succeq0$ 或者 $\nabla f(x)\neq0$ 且矩阵
\[
H(x)=
\begin{bmatrix}
\nabla^2 f(x) & \nabla f(x)\\
\nabla f(x)^T & 0
\end{bmatrix}
\]
只有一个负特征值。对任意 $y\neq0$，点 $x\in\dom f$ 满足式 (3.22) 的充要条件是 $H(x)$ 只有一个非正特征值。

提示：可以应用 (a) 的结论。下面的结论源自线性代数中的特征值交错定理，在本题的证明中也可能有用：如果矩阵 $B\in\mathbb{S}^n$，向量 $a\in\mathbb{R}^n$，那么有
\[
\lambda_n\left(
\begin{bmatrix}
B & a\\
a^T & 0
\end{bmatrix}
\right)\geq \lambda_n(B).
\]

\noindent\textbf{3.45}\quad 利用 \S3.4.3 给出的判断拟凸性的一阶和二阶条件来证明函数 $f(x)=-x_1x_2$ 的拟凸性，定义域 $\dom f=\mathbb{R}_{++}^2$。

\noindent\textbf{3.46\quad 定义域为 $\mathbb{R}^n$ 的拟线性函数。}\quad $\mathbb{R}$ 上的拟线性函数（既是拟凸函数又是拟凹函数）是单调的，即单调非减或者单调非增。在本题中，我们考虑上述结论扩展至 $\mathbb{R}^n$ 上的情形。

设函数 $f:\mathbb{R}^n\to\mathbb{R}$ 是连续的拟线性函数，其定义域为 $\dom f=\mathbb{R}^n$。证明函数可以表示为 $f(x)=g(a^Tx)$，其中函数 $g:\mathbb{R}\to\mathbb{R}$ 是单调函数，$a\in\mathbb{R}^n$。换言之，定义在 $\mathbb{R}^n$ 上的拟线性函数一定是某个线性函数的单调函数。（反过来同样正确。）

\noindent\textbf{对数-凹函数及对数-凸函数}

\noindent\textbf{3.47}\quad 设函数 $f:\mathbb{R}^n\to\mathbb{R}$ 可微，定义域 $\dom f$ 是凸集，对任意 $x\in\dom f$，$f(x)>0$。证明函数 $f$ 是对数-凹函数的充要条件是，对任意 $x,y\in\dom f$，下式成立
\[
\frac{f(y)}{f(x)}\leq \exp\left(\frac{\nabla f(x)^T(y-x)}{f(x)}\right).
\]

\noindent\textbf{3.48}\quad 证明如果函数 $f:\mathbb{R}^n\to\mathbb{R}$ 是对数-凹函数且 $a\geq0$，那么函数 $g=f-a$ 是对数-凹函数，其定义域为 $\dom g=\{x\in\dom f\mid f(x)>a\}$。

\noindent\textbf{3.49}\quad 证明下列函数是对数-凹函数。

(a) Logistic 函数
\[
f(x)=\frac{e^x}{1+e^x},
\]
其定义域为 $\dom f=\mathbb{R}$。

(b) 调和平均函数
\[
f(x)=\frac{1}{1/x_1+\cdots+1/x_n},\qquad \dom f=\mathbb{R}_{++}^n.
\]

(c) 乘积与和之比
\[
f(x)=\frac{\prod_{i=1}^n x_i}{\sum_{i=1}^n x_i},\qquad \dom f=\mathbb{R}_{++}^n.
\]

(d) 行列式与迹之比
\[
f(X)=\frac{\det X}{\operatorname{tr}X},\qquad \dom f=\mathbb{S}_{++}^n.
\]

\noindent\textbf{3.50\quad 将多项式的系数表示为根的函数。}\quad 证明具有负实根的多项式的系数可以表示为根的对数-凹函数。换言之，由如下等式定义的函数 $a_i:\mathbb{R}^n\to\mathbb{R}$，
\[
s^n+a_1(\lambda)s^{n-1}+\cdots+a_{n-1}(\lambda)s+a_n(\lambda)=(s-\lambda_1)(s-\lambda_2)\cdots(s-\lambda_n),
\]
在 $-\mathbb{R}_{++}^n$ 上是对数-凹函数。

提示：函数
\[
S_k(x)=\sum_{1\leq i_1<i_2<\cdots<i_k\leq n}x_{i_1}x_{i_2}\cdots x_{i_k},
\]
其定义域为 $\dom S_k=\mathbb{R}_+^n$，$1\leq k\leq n$，被称为 $\mathbb{R}^n$ 上的 $k$ 次初等对称函数。可以证明 $S_k^{1/k}$ 是凹函数（见参考文献 [ML57]）。

\noindent\textbf{3.51}\quad [BL00，第 41 页] 设 $p$ 是定义在 $\mathbb{R}$ 上的多项式，其所有的根都是实数。证明 $p$ 在所有使得它大于零的区间上是对数-凹函数。

\noindent\textbf{3.52}\quad [MO79，\S3.E.2] 矩函数的对数-凸性。设函数 $f:\mathbb{R}\to\mathbb{R}$ 在 $\mathbb{R}_+\subset\dom f$ 上非负。对 $x\geq0$，定义
\[
\phi(x)=\int_0^\infty u^x f(u)\,du.
\]
证明函数 $\phi$ 是对数-凸函数。（如果 $x$ 是正整数，$f$ 是概率密度函数，那么 $\phi(x)$ 就是此分布的 $x$ 阶矩。）

利用上述结论来证明 Gamma 函数
\[
\Gamma(x)=\int_0^\infty u^{x-1}e^{-u}\,du,
\]
在 $x\geq1$ 时是对数-凸函数。

\noindent\textbf{3.53}\quad 设 $x$ 和 $y$ 是 $\mathbb{R}^n$ 上的独立随机向量，其概率密度函数分别为 $f$ 和 $g$，均为对数-凹函数。证明随机向量的和 $z=x+y$ 的概率密度函数亦是对数-凹函数。

\noindent\textbf{3.54\quad Gauss 累积分布函数的对数-凹性。}\quad Gauss 随机变量的累积分布函数
\[
f(x)=\frac{1}{\sqrt{2\pi}}\int_{-\infty}^x e^{-t^2/2}\,dt,
\]
是对数-凹函数。根据两个对数-凹函数的卷积仍然是对数-凹函数，可以很容易得出 $f$ 是对数-凹函数的结论。在本题中，采用另一套证明方法证明 $f$ 的对数-凹性。我们知道，函数 $f$ 是对数-凹函数的充要条件是，对任意 $x$ 有 $f''(x)f(x)\leq f'(x)^2$。

(a) 证明对 $x\geq0$ 有 $f''(x)f(x)\leq f'(x)^2$。当 $x<0$ 时的情况较为复杂，在后续步骤中证明。

(b) 证明对任意 $t$ 和 $x$ 有 $t^2/2\geq -x^2/2+xt$。

(c) 利用 (b) 证明 $e^{-t^2/2}\leq e^{x^2/2-xt}$。因此，对 $x<0$ 有
\[
\int_{-\infty}^x e^{-t^2/2}\,dt\leq e^{x^2/2}\int_{-\infty}^x e^{-xt}\,dt.
\]

(d) 利用 (c) 证明当 $x\leq0$ 时，$f''(x)f(x)\leq f'(x)^2$。

\noindent\textbf{3.55\quad 对数-凹概率密度函数的累积分布函数的对数-凹性。}\quad 在这道题中，我们对习题 3.54 中的结论进行扩展。设函数 $g(t)=\exp(-h(t))$ 是一可微的对数-凹的概率密度函数，令
\[
f(x)=\int_{-\infty}^x g(t)\,dt=\int_{-\infty}^x e^{-h(t)}\,dt
\]
为其累积分布函数。下面的步骤将证明函数 $f$ 是对数-凹函数，即对任意 $x$ 成立 $f''(x)f(x)\leq(f'(x))^2$。

(a) 利用函数 $h$ 表示函数 $f$ 的导数。证明如果 $h'(x)\geq0$，那么 $f''(x)f(x)\leq(f'(x))^2$。

(b) 设 $h'(x)<0$。利用不等式
\[
h(t)\geq h(x)+h'(x)(t-x)
\]
（由函数 $h$ 的凸性可以直接得到）来证明
\[
\int_{-\infty}^x e^{-h(t)}\,dt\leq \frac{e^{-h(x)}}{-h'(x)}.
\]
利用上述不等式证明如果 $h'(x)<0$，有 $f''(x)f(x)\leq(f'(x))^2$。

\noindent\textbf{3.56\quad 更多的对数-凹的密度函数。}\quad 证明下列密度函数是对数-凹的。

(a) [MO79，第 493 页] Gamma 密度函数
\[
f(x)=\frac{\alpha^\lambda}{\Gamma(\lambda)}x^{\lambda-1}e^{-\alpha x},
\]
其定义域为 $\dom f=\mathbb{R}_+$。参数 $\lambda$ 和 $\alpha$ 满足 $\lambda\geq1,\alpha>0$。

(b) [MO79，第 306 页] Dirichlet 密度函数
\[
f(x)=\frac{\Gamma(\mathbf{1}^T\lambda)}{\Gamma(\lambda_1)\cdots\Gamma(\lambda_{n+1})}x_1^{\lambda_1-1}\cdots x_n^{\lambda_n-1}\left(1-\sum_{i=1}^n x_i\right)^{\lambda_{n+1}-1},
\]
其定义域为 $\dom f=\{x\in\mathbb{R}_{++}^n\mid \mathbf{1}^Tx<1\}$。参数 $\lambda$ 满足 $\lambda\succeq1$。

\noindent\textbf{关于广义不等式的凸性}

\noindent\textbf{3.57}\quad 证明函数 $f(X)=X^{-1}$ 在 $\mathbb{S}_{++}^n$ 上是矩阵凸的。

\noindent\textbf{3.58\quad Schur 补。}\quad 设矩阵 $X\in\mathbb{S}^n$，将其表示成分块矩阵的形式
\[
X=
\begin{bmatrix}
A & B\\
B^T & C
\end{bmatrix},
\]
其中 $A\in\mathbb{S}^k$。矩阵 $X$ 的 Schur 补（关于矩阵 $A$）为 $S=C-B^TA^{-1}B$（参见 \S A.5.5）。证明 Schur 补若看成 $\mathbb{S}^n$ 映射到 $\mathbb{S}^{n-k}$ 上的函数，在 $\mathbb{S}_{++}^n$ 上是矩阵凸的。

\noindent\textbf{3.59\quad $K$-凸的二阶条件。}\quad 设 $K\subseteq\mathbb{R}^m$ 是一正常凸锥，相应的广义不等式为 $\preceq_K$。证明定义域为凸集的二次可微函数 $f:\mathbb{R}^n\to\mathbb{R}^m$ 是 $K$-凸函数的充要条件是，对任意 $x\in\dom f$，任意 $y\in\mathbb{R}^n$，下式成立
\[
\sum_{i,j=1}^n \frac{\partial^2 f(x)}{\partial x_i\partial x_j}y_iy_j\succeq_K0,
\]
即二阶导数是 $K$-非负双线性形式。（此时 $\partial^2 f/\partial x_i\partial x_j\in\mathbb{R}^m$，分量为 $\partial^2 f_k/\partial x_i\partial x_j$，$k=1,\ldots,m$；参见 \S A.4.1。）

\noindent\textbf{3.60\quad $K$-凸函数的下水平集和上境图。}\quad 设 $K\subseteq\mathbb{R}^m$ 是正常凸锥，相应广义不等式为 $\preceq_K$，设 $f:\mathbb{R}^n\to\mathbb{R}^m$。对 $\alpha\in\mathbb{R}^m$，函数 $f$ 的 $\alpha$-下水平集（关于 $\preceq_K$）定义为
\[
C_\alpha=\{x\in\mathbb{R}^n\mid f(x)\preceq_K\alpha\}.
\]
函数 $f$ 关于 $\preceq_K$ 的上境图定义为集合
\[
\epi_K f=\{(x,t)\in\mathbb{R}^{n+m}\mid f(x)\preceq_K t\}.
\]
证明下列结论：

(a) 如果函数 $f$ 是 $K$-凸的，则对任意 $\alpha$，函数的下水平集 $C_\alpha$ 是凸集。

(b) 函数 $f$ 是 $K$-凸的充要条件是 $\epi_K f$ 是凸集。

\end{document}
